% !TEX root = ../高等概率论.tex
% =====================================================================
%  积分：简单函数、有界函数、非负函数、可积函数的积分，Levi 单调收敛定理，线性性，与 Riemann 积分的关系，Fatou 引理，控制收敛定理
% =====================================================================
\lecture{Riemann 积分的局限，简单函数的积分，有界函数的积分，非负可测函数的积分，Levi 单调收敛定理，积分的线性性，可积函数，Lebesgue 积分与 Riemann 积分，Fatou 引理，控制收敛定理，含参变量的积分}
\label{lec:integration}

有了测度和可测函数，我们现在来定义积分。在概率论中，积分就是期望：随机变量 $X$ 的期望 $EX$ 将被定义为 $\int X\dd P$。初等概率论中，期望有两套公式，离散型的 $\sum_xxP(X=x)$ 和有密度的 $\int xf(x)\dd x$。第~\ref{lec:distributions} 讲的 Cantor 分布既没有跳跃也没有密度，这两套公式对它都不适用。更何况 $\Omega$ 是一个抽象的集合，它上面没有“区间”，也就谈不上 Riemann 和。所以我们需要一个对任意测度空间都有意义的积分，并且它要把求和与 Riemann 积分都包括进来。

我们先回忆 Riemann 积分的长处与短处。按照定义，Riemann 积分是 Riemann 和 $\sum_kf(\xi_k)(x_k-x_{k-1})$ 的极限。用这个定义证明线性性是很简单的：$f+g$ 的 Riemann 和就是 $f$ 的 Riemann 和加上 $g$ 的 Riemann 和，取极限即可。可是涉及极限交换的性质就困难得多。比如说，设 $q_1,q_2,\dots$ 是 $[0,1]$ 中有理数的一个排列，令
\[
f_n=\one_{\{q_1,\dots,q_n\}}.
\]
每个 $f_n$ 只有有限个间断点，是 Riemann 可积的，积分为 $0$；$f_n$ 单调递增地收敛到 Dirichlet 函数 $\one_{\Q\cap[0,1]}$，而后者根本不是 Riemann 可积的。所以对 Riemann 积分，即使是单调递增的非负函数列，极限与积分也不能交换。数学分析中的办法是加上“一致收敛”的条件，而在概率论中，我们遇到的极限几乎都不是一致的。

Lebesgue 积分的做法正好反过来。它的定义方式使得线性性不再一目了然，但是极限与积分的交换非常方便。这一讲的核心结论是 \term{Levi 单调收敛定理}：若 $0\leqslant f_n\uparrow f$，则 $\int f_n\dd\mu\uparrow\int f\dd\mu$。我们会看到，只要有了这个定理，线性性就可以通过简单函数逼近几乎自动地得到。

定义分四步完成：简单函数，有界函数，非负函数，一般的函数。每一步都证明同样的几条性质。这种“四步法”以后在证明关于积分的等式时还会反复使用。

\topic{本讲用到的前文结果}

\begin{enumerate}
\item 测度的基本性质与连续性（定理~\ref{thm:measures:basic}）：有限可加、单调、次可加；$A_n\uparrow A$ 时 $\mu(A_n)\uparrow\mu(A)$；$A_n\downarrow A$ 且 $\mu(A_1)<\infty$ 时 $\mu(A_n)\downarrow\mu(A)$。
\item $\sigma$-有限测度（定义~\ref{def:measures:finite}）：存在 $E_n\uparrow\Omega$，$\mu(E_n)<\infty$。
\item 可测函数的运算与极限（推论~\ref{cor:random-variables:arithmetic}、定理~\ref{thm:random-variables:limits}）：和、积、$f^+$、$f^-$、$\sup_n f_n$、$\lim_nf_n$ 都可测。
\item 简单函数逼近（定理~\ref{thm:random-variables:simple-approx}）：非负可测函数 $f$ 是递增的非负简单函数列 $\varphi_n$ 的逐点极限。
\item Lebesgue 测度：区间的测度等于长度，单点与可数集是零测集（命题~\ref{prop:extension:lebesgue}）。
\item 数学分析中的 Riemann 积分：Darboux 上和 $U(P)$、下和 $L(P)$，Riemann 可积当且仅当 $\sup_PL(P)=\inf_PU(P)$。
\end{enumerate}

\emph{本讲约定：}$(\Omega,\cF,\mu)$ 是测度空间，并且 $\mu$ 是 $\sigma$-有限的（概率测度当然是有限的，从而是 $\sigma$-有限的）。所说的函数都是 $\cF$-可测的。在 $[0,\infty]$ 中约定 $0\cdot\infty=0$。

\topic{简单函数的积分}

回忆第~\ref{lec:random-variables} 讲的定义~\ref{def:random-variables:simple}：简单函数是形如 $\sum_{i=1}^nc_i\one_{A_i}$（$A_i\in\cF$）的函数，也就是只取有限个值的可测函数。若简单函数 $\varphi$ 取的不同的值是 $y_1,\dots,y_m$，则
\[
\varphi=\sum_{j=1}^my_j\one_{\{\varphi=y_j\}},
\]
其中 $\{\varphi=y_j\}\in\cF$ 两两不交，并为 $\Omega$。这称为 $\varphi$ 的\term{标准表示}（canonical representation）。

\begin{definition}[简单函数的积分]\label{def:integration:simple}
设 $\varphi$ 是简单函数，并且 $\mu(\{\varphi\neq0\})<\infty$。定义
\[
\int\varphi\dd\mu=\sum_{y\neq0}y\,\mu(\{\varphi=y\}),
\]
其中 $y$ 取遍 $\varphi$ 取到的一切非零值（只有有限个）。
\end{definition}

由于 $\{\varphi=y\}\subseteq\{\varphi\neq0\}$（$y\neq0$），和式中每一项都是有限的实数。我们要求 $\mu(\{\varphi\neq0\})<\infty$，是为了避免出现 $\infty-\infty$：比如在 $(\R,\cR,\lambda)$ 上，$\one_{[0,\infty)}-\one_{(-\infty,0)}$ 的积分就没有合理的定义。

在计算中，简单函数常常不是以标准表示给出的，比如 $\varphi=2\cdot\one_{[0,1]}+2\cdot\one_{(1,3]}$，同一个值 $2$ 分成了两块。下面的引理说，只要各块两两不交，按哪种表示计算都一样。

\begin{lemma}\label{lem:integration:representation}
设 $\varphi=\sum_{i=1}^na_i\one_{A_i}$，其中 $a_i\in\R$，$A_i\in\cF$ 两两不交，并且 $a_i\neq0$ 时 $\mu(A_i)<\infty$。则 $\mu(\{\varphi\neq0\})<\infty$，且
\[
\int\varphi\dd\mu=\sum_{i=1}^na_i\mu(A_i).
\]
\end{lemma}
\begin{proof}
由于 $A_i$ 两两不交，每个 $\omega$ 至多属于一个 $A_i$，所以 $\omega\in A_i$ 时 $\varphi(\omega)=a_i$，$\omega\notin\bigcup_iA_i$ 时 $\varphi(\omega)=0$。于是对 $\varphi$ 取到的每个非零值 $y$，
\[
\{\varphi=y\}=\bigcup_{i:\,a_i=y}A_i\quad(\text{不交并}).
\]
事实上，若 $\omega\in A_i$ 且 $a_i=y$，则 $\varphi(\omega)=y$；反之若 $\varphi(\omega)=y\neq0$，则 $\omega$ 属于某个 $A_i$，并且 $a_i=\varphi(\omega)=y$。特别地 $\{\varphi\neq0\}=\bigcup_{i:a_i\neq0}A_i$ 的测度有限。由有限可加性，
\[
\int\varphi\dd\mu=\sum_{y\neq0}y\sum_{i:\,a_i=y}\mu(A_i)=\sum_{i:\,a_i\neq0}a_i\mu(A_i)=\sum_{i=1}^na_i\mu(A_i),
\]
最后一步是因为 $a_i=0$ 的项等于 $0\cdot\mu(A_i)=0$（即使 $\mu(A_i)=\infty$）。
\end{proof}

记 $\cE_0$ 为满足 $\mu(\{\varphi\neq0\})<\infty$ 的简单函数的全体。若 $\varphi,\psi\in\cE_0$，$a\in\R$，则 $a\varphi$，$\varphi+\psi$，$\abs\varphi$ 都属于 $\cE_0$，因为 $\{\varphi+\psi\neq0\}\subseteq\{\varphi\neq0\}\cup\{\psi\neq0\}$。

\begin{proposition}\label{prop:integration:simple}
设 $\varphi,\psi\in\cE_0$。则
\begin{enumerate}
\item 若 $\varphi\geqslant0$ a.e.，则 $\int\varphi\dd\mu\geqslant0$；
\item 对任意 $a\in\R$，$\int a\varphi\dd\mu=a\int\varphi\dd\mu$；
\item $\int(\varphi+\psi)\dd\mu=\int\varphi\dd\mu+\int\psi\dd\mu$。
\end{enumerate}
\end{proposition}
\begin{proof}
(1) 若 $y<0$ 是 $\varphi$ 取到的值，则 $\{\varphi=y\}\subseteq\{\varphi<0\}$ 是零测集，这一项为 $0$。其余各项 $y\mu(\{\varphi=y\})$ 都 $\geqslant0$。

(2) $a=0$ 时两边都是 $0$。$a\neq0$ 时，$a\varphi$ 取的非零值恰好是 $ay$（$y$ 取遍 $\varphi$ 的非零值），且 $\{a\varphi=ay\}=\{\varphi=y\}$，所以 $\int a\varphi\dd\mu=\sum_{y\neq0}ay\,\mu(\{\varphi=y\})=a\int\varphi\dd\mu$。

(3) 设 $\varphi=\sum_{i=1}^ma_i\one_{A_i}$，$\psi=\sum_{j=1}^nb_j\one_{B_j}$ 是只保留非零值的标准表示，即 $a_i$ 是 $\varphi$ 的全部非零值，$A_i=\{\varphi=a_i\}$；$b_j$，$B_j$ 同样。为了让两组集合“覆盖同一块地方”，令
\[
A_0=\bigcup_jB_j\setminus\bigcup_iA_i,\qquad B_0=\bigcup_iA_i\setminus\bigcup_jB_j,\qquad a_0=b_0=0.
\]
则 $\varphi=\sum_{i=0}^ma_i\one_{A_i}$，$\psi=\sum_{j=0}^nb_j\one_{B_j}$，$A_0,\dots,A_m$ 两两不交，$B_0,\dots,B_n$ 两两不交，并且二者的并都等于 $U=\bigcup_iA_i\cup\bigcup_jB_j$，$\mu(U)<\infty$。于是 $A_i=\bigcup_{j=0}^n(A_i\cap B_j)$，$B_j=\bigcup_{i=0}^m(A_i\cap B_j)$，都是不交并，并且
\[
\varphi+\psi=\sum_{i=0}^m\sum_{j=0}^n(a_i+b_j)\one_{A_i\cap B_j},
\]
其中 $A_i\cap B_j$ 两两不交，测度有限。由引理~\ref{lem:integration:representation} 与有限可加性，
\[
\begin{aligned}
\int(\varphi+\psi)\dd\mu&=\sum_{i,j}(a_i+b_j)\mu(A_i\cap B_j)=\sum_ia_i\sum_j\mu(A_i\cap B_j)+\sum_jb_j\sum_i\mu(A_i\cap B_j)\\
&=\sum_ia_i\mu(A_i)+\sum_jb_j\mu(B_j),
\end{aligned}
\]
这就是 $\int\varphi\dd\mu+\int\psi\dd\mu$。
\end{proof}

下面三条性质只用到 (1)(2)(3)，以后每推广一次积分，只要验证了 (1)(2)(3)，它们就自动成立。为了不重复同样的论证，我们把它写成一个一般的引理。

\begin{lemma}\label{lem:integration:consequences}
设 $\cK$ 是一族实值函数，对加法、数乘和取绝对值封闭，$I:\cK\to\R$ 满足：(1) $f\geqslant0$ a.e. 蕴含 $I(f)\geqslant0$；(2) $I(af)=aI(f)$；(3) $I(f+g)=I(f)+I(g)$。则对 $f,g\in\cK$：
\begin{enumerate}
\setcounter{enumi}{3}
\item 若 $f\leqslant g$ a.e.，则 $I(f)\leqslant I(g)$；
\item 若 $f=g$ a.e.，则 $I(f)=I(g)$；
\item $\abs{I(f)}\leqslant I(\abs f)$。
\end{enumerate}
\end{lemma}
\begin{proof}
(4) $g-f=g+(-1)f\in\cK$ 且 $g-f\geqslant0$ a.e.，由 (3)(2)(1)，$I(g)=I(f)+I(g-f)\geqslant I(f)$。(5) $f=g$ a.e. 蕴含 $f\leqslant g$ a.e. 与 $g\leqslant f$ a.e.，用两次 (4)。(6) 由 $f\leqslant\abs f$ 与 (4) 得 $I(f)\leqslant I(\abs f)$；由 $-f\leqslant\abs f$，(4) 与 (2) 得 $-I(f)\leqslant I(\abs f)$。
\end{proof}

特别地，$\cE_0$ 上的积分满足 (4)(5)(6)。

\topic{有界函数的积分}

第二步，设 $E\in\cF$，$\mu(E)<\infty$，$f$ 是有界可测函数，并且在 $E^c$ 上 $f=0$。若 $\varphi,\psi\in\cE_0$ 满足 $\varphi\leqslant f\leqslant\psi$，我们当然希望 $\int\varphi\dd\mu\leqslant\int f\dd\mu\leqslant\int\psi\dd\mu$。这启发我们用下面的方式定义 $\int f\dd\mu$。以下在本节中，$\varphi,\psi$ 总是指在 $E^c$ 上为零的简单函数（它们自动属于 $\cE_0$）。

\begin{theorem}\label{thm:integration:bounded}
在上述假设下，
\[
\sup_{\varphi\leqslant f}\int\varphi\dd\mu=\inf_{\psi\geqslant f}\int\psi\dd\mu .
\]
这个公共值定义为 $\int f\dd\mu$。
\end{theorem}
\begin{proof}
设 $\abs f\leqslant M$。$\varphi=-M\one_E$ 与 $\psi=M\one_E$ 分别满足 $\varphi\leqslant f\leqslant\psi$，所以上确界与下确界都是对非空集合取的。若 $\varphi\leqslant f\leqslant\psi$，则 $\varphi\leqslant\psi$，由引理~\ref{lem:integration:consequences}(4)，$\int\varphi\dd\mu\leqslant\int\psi\dd\mu$。先对 $\varphi$ 取上确界、再对 $\psi$ 取下确界，得到“$\leqslant$”。

反过来，固定正整数 $n$，把值域 $[-M,M]$ 等分成 $2n$ 段，令
\[
E_k=\Bigl\{\omega\in E:\ \frac{(k-1)M}n<f(\omega)\leqslant\frac{kM}n\Bigr\},\qquad -n\leqslant k\leqslant n,
\]
其中 $k=-n$ 时的集合就是 $\{\omega\in E: f(\omega)=-M\}$。由于 $f$ 可测，$E_k\in\cF$；$E_{-n},\dots,E_n$ 两两不交，并且它们的并是 $E$，因为 $[-M,M]$ 是 $\{-M\}$ 与 $2n$ 个区间 $((k-1)M/n,kM/n]$（$-n<k\leqslant n$）的不交并。令
\[
\psi_n=\sum_{k=-n}^n\frac{kM}n\one_{E_k},\qquad \varphi_n=\sum_{k=-n}^n\frac{(k-1)M}n\one_{E_k}.
\]
则 $\varphi_n\leqslant f\leqslant\psi_n$，并且 $\psi_n-\varphi_n=\frac Mn\one_E$。由命题~\ref{prop:integration:simple}，
\[
\sup_{\varphi\leqslant f}\int\varphi\dd\mu\geqslant\int\varphi_n\dd\mu=\int\psi_n\dd\mu-\frac Mn\mu(E)\geqslant\inf_{\psi\geqslant f}\int\psi\dd\mu-\frac Mn\mu(E).
\]
由于 $\mu(E)<\infty$，令 $n\to\infty$ 即得“$\geqslant$”。
\end{proof}

\begin{remark}[与 Darboux 和比较]
这个定义在形式上和 Riemann 积分的 Darboux 上、下和完全一样：从下面和上面用“简单”的函数夹住 $f$，看上、下两个极限是否相等。区别在于什么样的函数算“简单”。在 Riemann 积分中，夹住 $f$ 的是阶梯函数，它们只能在区间上取常数，这相当于分割\emph{定义域}。对 Dirichlet 函数 $\one_{\Q\cap[0,1]}$，任何一个小区间里都既有有理数又有无理数，阶梯函数从上面夹住它只能取 $1$，从下面夹住它只能取 $0$，上、下和永远相差 $1$。而在这里，$\varphi,\psi$ 可以在任意可测集上取常数，于是我们可以分割\emph{值域}：定理的证明中，$\varphi_n$ 与 $\psi_n$ 的差恒为 $M/n$，与 $f$ 多么不连续毫无关系。这样做的代价是，集合 $E_k=\{(k-1)M/n<f\leqslant kM/n\}$ 必须有测度，这正是我们要求 $f$ 可测的原因。所以在 Lebesgue 的理论中，有界可测函数在测度有限的集合上总是可积的。
\end{remark}

我们指出这一步定义的两个一致性。第一，若 $f$ 本身属于 $\cE_0$ 且在 $E^c$ 上为零，取 $\varphi=\psi=f$，可见新的定义与定义~\ref{def:integration:simple} 一致。第二，定义与 $E$ 的选取无关：设 $f$ 在 $E^c$ 与 $E'^c$ 上都为零，$\mu(E),\mu(E')<\infty$，则 $f$ 在 $(E\cap E')^c$ 上也为零。若 $\varphi\leqslant f$ 在 $E^c$ 上为零，则在 $E\setminus E'$ 上 $\varphi\leqslant f=0$，所以 $\varphi\one_{E\cap E'}\geqslant\varphi$，并且 $\varphi\one_{E\cap E'}\leqslant f$（在 $E\cap E'$ 上它等于 $\varphi$，在别处它等于 $0=f$）。由引理~\ref{lem:integration:consequences}(4)，用 $E$ 算出的上确界不超过用 $E\cap E'$ 算出的；反过来的不等式是显然的，因为在 $(E\cap E')^c$ 上为零的 $\varphi$ 也在 $E^c$ 上为零。所以用 $E$ 与用 $E\cap E'$ 定义的积分相同，同理用 $E'$ 与用 $E\cap E'$ 定义的也相同。

\begin{proposition}\label{prop:integration:bounded}
设 $\mu(E)<\infty$，$f,g$ 是有界可测函数，在 $E^c$ 上为零。则
\begin{enumerate}
\item 若 $f\geqslant0$ a.e.，则 $\int f\dd\mu\geqslant0$；
\item 对任意 $a\in\R$，$\int af\dd\mu=a\int f\dd\mu$；
\item $\int(f+g)\dd\mu=\int f\dd\mu+\int g\dd\mu$；
\end{enumerate}
从而由引理~\ref{lem:integration:consequences}，(4)(5)(6) 也成立。
\end{proposition}
\begin{proof}
(1) 设 $N=\{f<0\}$，$\mu(N)=0$。简单函数 $\varphi=-M\one_{E\cap N}$（$M$ 是 $\abs f$ 的上界）满足 $\varphi\leqslant f$，并且 $\int\varphi\dd\mu=-M\mu(E\cap N)=0$。所以 $\int f\dd\mu\geqslant0$。

(2) $a=0$ 显然。设 $a>0$。$a\varphi\leqslant af$ 当且仅当 $\varphi\leqslant f$，而 $\varphi\mapsto a\varphi$ 是在 $E^c$ 上为零的简单函数全体到自身的一一对应，所以
\[
\int af\dd\mu=\sup_{\varphi\leqslant f}\int a\varphi\dd\mu=a\sup_{\varphi\leqslant f}\int\varphi\dd\mu=a\int f\dd\mu .
\]
设 $a<0$。此时 $a\psi\leqslant af$ 当且仅当 $\psi\geqslant f$，所以
\[
\int af\dd\mu=\sup_{\psi\geqslant f}\int a\psi\dd\mu=a\inf_{\psi\geqslant f}\int\psi\dd\mu=a\int f\dd\mu,
\]
这里用了 $a<0$ 时 $\sup(a\,\cdot)=a\inf(\cdot)$。

(3) $f+g$ 有界，在 $E^c$ 上为零。若 $\psi_1\geqslant f$，$\psi_2\geqslant g$，则 $\psi_1+\psi_2\geqslant f+g$，由定理~\ref{thm:integration:bounded} 中的下确界表示与简单函数积分的线性性，
\[
\int(f+g)\dd\mu\leqslant\int(\psi_1+\psi_2)\dd\mu=\int\psi_1\dd\mu+\int\psi_2\dd\mu .
\]
对 $\psi_1$、$\psi_2$ 分别取下确界得 $\int(f+g)\dd\mu\leqslant\int f\dd\mu+\int g\dd\mu$。把这个不等式用于 $-f,-g$，再用 (2)，得 $-\int(f+g)\dd\mu\leqslant-\int f\dd\mu-\int g\dd\mu$，即反向不等式。
\end{proof}

对 $A\in\cF$，记 $\int_Af\dd\mu=\int f\one_A\dd\mu$，称为 $f$ 在 $A$ 上的积分。

\topic{非负函数的积分}

第三步，设 $f:\Omega\to[0,\infty]$ 可测。

\begin{definition}[非负函数的积分]\label{def:integration:nonnegative}
\[
\int f\dd\mu=\sup\Bigl\{\int h\dd\mu:\ 0\leqslant h\leqslant f,\ h\ \text{有界可测},\ \mu(\{h>0\})<\infty\Bigr\}\in[0,\infty].
\]
\end{definition}

这里的每个 $h$ 在 $\{h>0\}^c$ 上为零，而 $\{h>0\}$ 的测度有限，所以 $\int h\dd\mu$ 已经由第二步定义好了。这个定义的好处是它显然有意义，不需要再证明什么“上确界等于下确界”。我们仍然先说明它与第二步一致：若 $f\geqslant0$ 有界，并且在测度有限的集合 $E$ 之外为零，那么 $f$ 本身就是定义中的一个 $h$，所以新的积分不小于旧的；另一方面，定义中的每个 $h$ 满足 $0\leqslant h\leqslant f$，都在 $E^c$ 上为零，由命题~\ref{prop:integration:bounded}(4)，$\int h\dd\mu$ 不超过旧的积分。所以二者相等。

由定义立即可以得到下面几条性质。

\begin{proposition}\label{prop:integration:nonnegative-basic}
设 $f,g:\Omega\to[0,\infty]$ 可测。
\begin{enumerate}
\item 若 $g\leqslant f$，则 $\int g\dd\mu\leqslant\int f\dd\mu$。
\item 对 $0\leqslant a<\infty$，$\int af\dd\mu=a\int f\dd\mu$。
\item 若 $g\leqslant f$ a.e.，则 $\int g\dd\mu\leqslant\int f\dd\mu$；若 $f=g$ a.e.，则 $\int f\dd\mu=\int g\dd\mu$。
\end{enumerate}
\end{proposition}
\begin{proof}
(1) 定义 $\int g\dd\mu$ 的那一族 $h$ 包含于定义 $\int f\dd\mu$ 的那一族之中，前者的上确界不超过后者。

(2) $a=0$ 时两边都是 $0$（$0\cdot f$ 恒为零，而零函数的积分为 $0$）。$a>0$ 时，$h\mapsto ah$ 是定义 $\int f$ 的那一族到定义 $\int af$ 的那一族的一一对应，且 $\int ah\dd\mu=a\int h\dd\mu$（命题~\ref{prop:integration:bounded}(2)）。

(3) 设 $N$ 是零测集，在 $N^c$ 上 $g\leqslant f$。任取定义 $\int g$ 的一个 $h$，令 $h'=h\one_{N^c}$。在 $N^c$ 上 $h'=h\leqslant g\leqslant f$，在 $N$ 上 $h'=0\leqslant f$，所以 $h'$ 属于定义 $\int f$ 的那一族；又 $h'=h$ a.e.，由命题~\ref{prop:integration:bounded}(5)，$\int h\dd\mu=\int h'\dd\mu\leqslant\int f\dd\mu$。对 $h$ 取上确界得第一个结论。第二个结论由第一个结论用两次得到。
\end{proof}

下面的引理用到了 $\mu$ 的 $\sigma$-有限性，它给出了计算积分的一个办法：把 $f$ 在高度 $n$ 处截断、在测度有限的集合 $E_n$ 之外砍掉，得到的是第二步中的有界函数。

\begin{lemma}\label{lem:integration:truncation}
设 $E_n\in\cF$，$E_n\uparrow\Omega$，$\mu(E_n)<\infty$（由 $\sigma$-有限性，这样的 $E_n$ 存在）。记 $a\wedge b=\min\{a,b\}$。则对可测的 $f:\Omega\to[0,\infty]$，
\[
\int_{E_n}(f\wedge n)\dd\mu\ \uparrow\ \int f\dd\mu\qquad(n\to\infty).
\]
\end{lemma}
\begin{proof}
$h_n=(f\wedge n)\one_{E_n}$ 有界，在 $E_n^c$ 上为零，并且 $h_n\leqslant h_{n+1}$，由命题~\ref{prop:integration:bounded}(4)，左边关于 $n$ 单调不减。$h_n$ 是定义~\ref{def:integration:nonnegative} 中的一个 $h$，所以左边每一项不超过 $\int f\dd\mu$。

为了证明极限等于 $\int f\dd\mu$，任取定义中的一个 $h$：$0\leqslant h\leqslant f$，$h\leqslant M$，$\mu(\{h>0\})<\infty$。当 $n\geqslant M$ 时，$h\leqslant f\wedge n$，所以 $h\one_{E_n}\leqslant h_n$。由命题~\ref{prop:integration:bounded}(3)(4)，
\[
\int h_n\dd\mu\geqslant\int h\one_{E_n}\dd\mu=\int h\dd\mu-\int h\one_{E_n^c}\dd\mu .
\]
而 $0\leqslant h\one_{E_n^c}\leqslant M\one_{E_n^c\cap\{h>0\}}$，所以 $0\leqslant\int h\one_{E_n^c}\dd\mu\leqslant M\mu(E_n^c\cap\{h>0\})$。由于 $E_n^c\cap\{h>0\}\downarrow\varnothing$，且 $\mu(\{h>0\})<\infty$，由测度的从上连续，$\mu(E_n^c\cap\{h>0\})\to0$。于是 $\lim_n\int h_n\dd\mu\geqslant\int h\dd\mu$。对 $h$ 取上确界即得。
\end{proof}

作为第一个应用，我们算出非负简单函数的积分，这里不再要求 $\{\varphi>0\}$ 的测度有限。

\begin{corollary}\label{cor:integration:nonnegative-simple}
设 $\varphi\geqslant0$ 是简单函数，取的不同的值为 $y_1,\dots,y_m$。则 $\int\varphi\dd\mu=\sum_{j=1}^my_j\mu(\{\varphi=y_j\})$（约定 $0\cdot\infty=0$）。特别地，$\int\one_A\dd\mu=\mu(A)$ 对一切 $A\in\cF$ 成立。
\end{corollary}
\begin{proof}
当 $n\geqslant\max_jy_j$ 时 $\varphi\wedge n=\varphi$，于是 $(\varphi\wedge n)\one_{E_n}=\sum_jy_j\one_{\{\varphi=y_j\}\cap E_n}$ 属于 $\cE_0$，由引理~\ref{lem:integration:representation}，它的积分是 $\sum_jy_j\mu(\{\varphi=y_j\}\cap E_n)$。由测度的从下连续，$\mu(\{\varphi=y_j\}\cap E_n)\uparrow\mu(\{\varphi=y_j\})$。再由引理~\ref{lem:integration:truncation} 即得。
\end{proof}

\topic{Levi 单调收敛定理}

现在证明本讲的核心定理。

\begin{remark}[证明思路]
$\int f_n\leqslant\int f$ 是显然的，难的是反向。按定义，$\int f$ 是 $f$ 下方的有界函数 $h$ 的积分的上确界，所以只需证明 $\lim_n\int f_n\geqslant\int h$。我们在集合 $A_n=\{f_n\geqslant ch\}$（$0<c<1$）上用 $f_n\geqslant ch$，而 $A_n$ 逐渐长满全空间，剩下的部分的测度由测度的连续性趋于零。定理之后的注记会把这个思路再说一遍。
\end{remark}

\begin{theorem}[Levi 单调收敛定理]\label{thm:integration:mct}
设 $f_n:\Omega\to[0,\infty]$ 可测，$f_1\leqslant f_2\leqslant\cdots$，$f_n\uparrow f$（逐点）。则
\[
\int f_n\dd\mu\ \uparrow\ \int f\dd\mu .
\]
\end{theorem}
\begin{proof}
由定理~\ref{thm:random-variables:limits}，$f=\sup_nf_n$ 是可测的。由命题~\ref{prop:integration:nonnegative-basic}(1)，$\int f_n\dd\mu$ 单调不减，且每一项不超过 $\int f\dd\mu$。记 $L=\lim_n\int f_n\dd\mu$，则 $L\leqslant\int f\dd\mu$。

为了证明 $L\geqslant\int f\dd\mu$，只需对定义~\ref{def:integration:nonnegative} 中的每个 $h$ 证明 $L\geqslant\int h\dd\mu$。设 $0\leqslant h\leqslant f$，$h\leqslant M$，$S=\{h>0\}$，$\mu(S)<\infty$。固定 $c\in(0,1)$，令
\[
A_n=\{f_n\geqslant ch\}.
\]
$A_n\in\cF$：它的补集 $\{f_n<ch\}=\bigcup_{q\in\Q}(\{f_n<q\}\cap\{q<ch\})$ 是可数个事件的并。由于 $f_n$ 单调不减，$A_n$ 单调递增。我们断言 $A_n\uparrow\Omega$：若 $h(\omega)=0$，则 $\omega\in A_1$；若 $h(\omega)>0$，则 $f(\omega)\geqslant h(\omega)>ch(\omega)$，而 $f_n(\omega)\uparrow f(\omega)$，所以 $n$ 充分大时 $f_n(\omega)>ch(\omega)$。

函数 $ch\one_{A_n}$ 有界，在 $S^c$ 上为零，并且 $ch\one_{A_n}\leqslant f_n$（在 $A_n$ 上由 $A_n$ 的定义，在 $A_n^c$ 上左边为零）。所以它是定义 $\int f_n\dd\mu$ 的一个 $h$，从而由命题~\ref{prop:integration:bounded}(2)(3)，
\[
\int f_n\dd\mu\geqslant\int ch\one_{A_n}\dd\mu=c\Bigl(\int h\dd\mu-\int h\one_{A_n^c}\dd\mu\Bigr)\geqslant c\Bigl(\int h\dd\mu-M\mu(S\cap A_n^c)\Bigr).
\]
由于 $S\cap A_n^c\downarrow\varnothing$，$\mu(S)<\infty$，由从上连续，$\mu(S\cap A_n^c)\to0$。令 $n\to\infty$ 得 $L\geqslant c\int h\dd\mu$。再令 $c\uparrow1$ 得 $L\geqslant\int h\dd\mu$。
\end{proof}

\begin{remark}
这是一个值得学习的证明，我们把它的思路说清楚。要证明的是 $f_n$ 的积分最终不会比 $f$ 的积分小多少。$f$ 的积分是由 $f$ 下方的有界函数 $h$ 逼近的，所以只要说明 $f_n$ 最终能“托住” $h$ 的几乎全部。$f_n$ 未必处处不小于 $h$，但在集合 $A_n$ 上它不小于 $ch$，而 $A_n$ 逐渐长满整个空间。于是差别只在 $A_n^c$ 上，它的测度趋于零。取 $c<1$ 而不是 $c=1$，是因为 $f_n$ 可能永远严格小于 $f=h$，这时 $\{f_n\geqslant h\}$ 可能是空集。

证明中关键的一步是 $\mu(S\cap A_n^c)\to0$，它用到的是测度的从上连续，也就是可数可加性。这正是 Riemann 积分所缺少的：在本讲开头的例子中，$f_n=\one_{\{q_1,\dots,q_n\}}$ 的极限是有理数集的示性函数，而“$[0,1]$ 中有理数集的长度是 $0$”这个结论，需要把可数个单点的长度加起来。Jordan 测度只有有限可加性，做不到这一点。
\end{remark}

\topic{积分的线性性}

从定义~\ref{def:integration:nonnegative} 出发，线性性的一半是容易的：若 $h$ 是定义 $\int f$ 的一个函数，$k$ 是定义 $\int g$ 的一个函数，则 $h+k$ 有界，$\{h+k>0\}\subseteq\{h>0\}\cup\{k>0\}$ 的测度有限，并且 $h+k\leqslant f+g$，所以
\[
\int(f+g)\dd\mu\geqslant\int(h+k)\dd\mu=\int h\dd\mu+\int k\dd\mu .
\]
对 $h,k$ 取上确界得 $\int(f+g)\dd\mu\geqslant\int f\dd\mu+\int g\dd\mu$。反向的不等式就不容易了：给定 $f+g$ 下方的一个 $h$，我们需要把它拆成 $f$ 下方的一块与 $g$ 下方的一块，这并不方便。这时单调收敛定理就派上了用场：用简单函数从下面逼近 $f$ 和 $g$，对简单函数线性性已经成立，再取极限。

\begin{theorem}[非负函数积分的线性性]\label{thm:integration:linear-nonnegative}
设 $f,g:\Omega\to[0,\infty]$ 可测，$a,b\in[0,\infty)$。则
\[
\int(af+bg)\dd\mu=a\int f\dd\mu+b\int g\dd\mu .
\]
\end{theorem}
\begin{proof}
由命题~\ref{prop:integration:nonnegative-basic}(2)，只需证明 $\int(f+g)\dd\mu=\int f\dd\mu+\int g\dd\mu$。

\emph{第一步：非负简单函数。}设 $\varphi,\psi\geqslant0$ 是简单函数（不要求 $\{\varphi>0\}$ 的测度有限）。取 $E_n\uparrow\Omega$，$\mu(E_n)<\infty$。则 $\varphi\one_{E_n}$，$\psi\one_{E_n}$ 属于 $\cE_0$，由命题~\ref{prop:integration:simple}(3)，
\[
\int(\varphi+\psi)\one_{E_n}\dd\mu=\int\varphi\one_{E_n}\dd\mu+\int\psi\one_{E_n}\dd\mu .
\]
（由前两步的一致性，这里的积分按哪一步的定义理解都一样。）当 $n\to\infty$ 时，$\varphi\one_{E_n}\uparrow\varphi$，$\psi\one_{E_n}\uparrow\psi$，$(\varphi+\psi)\one_{E_n}\uparrow\varphi+\psi$，由单调收敛定理，三个积分分别收敛到 $\int\varphi$、$\int\psi$、$\int(\varphi+\psi)$。所以 $\int(\varphi+\psi)\dd\mu=\int\varphi\dd\mu+\int\psi\dd\mu$。

\emph{第二步：一般情形。}由定理~\ref{thm:random-variables:simple-approx}，存在非负简单函数 $\varphi_n\uparrow f$，$\psi_n\uparrow g$。于是 $\varphi_n+\psi_n$ 是非负简单函数，并且 $\varphi_n+\psi_n\uparrow f+g$（在 $[0,\infty]$ 中，单调数列的和的极限等于极限的和）。由单调收敛定理与第一步，
\[
\int(f+g)\dd\mu=\lim_n\int(\varphi_n+\psi_n)\dd\mu=\lim_n\Bigl(\int\varphi_n\dd\mu+\int\psi_n\dd\mu\Bigr)=\int f\dd\mu+\int g\dd\mu .
\]
\end{proof}

\begin{remark}
我们看到，单调收敛定理是 Lebesgue 积分理论的核心：积分的定义本身只提供了“从下面逼近”这一个信息，而单调收敛定理说，\emph{任何}一列从下面单调逼近 $f$ 的函数都给出同一个积分。有了它，线性性、逐项积分以及下面的一系列性质都顺理成章。上面的证明也展示了今后证明积分等式的标准套路：先对示性函数验证，由线性性推到简单函数，由单调收敛定理推到非负可测函数，最后通过 $f=f^+-f^-$ 推到一般的可积函数。
\end{remark}

下面是线性性与单调收敛定理的几个直接推论。

\begin{corollary}[逐项积分]\label{cor:integration:series}
设 $f_n:\Omega\to[0,\infty]$ 可测。则 $\int\sum_{n=1}^\infty f_n\dd\mu=\sum_{n=1}^\infty\int f_n\dd\mu$。
\end{corollary}
\begin{proof}
部分和 $g_N=\sum_{n\leqslant N}f_n$ 非负可测，单调递增地收敛到 $\sum_nf_n$。由定理~\ref{thm:integration:linear-nonnegative}，$\int g_N\dd\mu=\sum_{n\leqslant N}\int f_n\dd\mu$。由单调收敛定理，令 $N\to\infty$ 即得。
\end{proof}

请同学们注意，这里对非负函数不需要任何收敛性的假设，这与数学分析中逐项积分需要一致收敛形成鲜明的对比。

\begin{corollary}[不定积分是测度]\label{cor:integration:indefinite}
设 $f:\Omega\to[0,\infty]$ 可测。则 $\nu(A)=\int_Af\dd\mu$（$A\in\cF$）是 $(\Omega,\cF)$ 上的测度。
\end{corollary}
\begin{proof}
$\nu(\varnothing)=\int0\dd\mu=0$。设 $A_n\in\cF$ 两两不交，$A=\bigcup_nA_n$。则 $f\one_A=\sum_nf\one_{A_n}$（每一点至多有一项非零），由推论~\ref{cor:integration:series}，$\nu(A)=\sum_n\int f\one_{A_n}\dd\mu=\sum_n\nu(A_n)$。
\end{proof}

\begin{proposition}\label{prop:integration:markov}
设 $f:\Omega\to[0,\infty]$ 可测。
\begin{enumerate}
\item （Markov 不等式）对 $c>0$，$\mu(\{f\geqslant c\})\leqslant\frac1c\int f\dd\mu$。
\item $\int f\dd\mu=0$ 当且仅当 $f=0$ a.e.
\item 若 $\int f\dd\mu<\infty$，则 $f<\infty$ a.e.
\end{enumerate}
\end{proposition}
\begin{proof}
(1) $c\one_{\{f\geqslant c\}}\leqslant f$，由命题~\ref{prop:integration:nonnegative-basic}(1)(2) 与推论~\ref{cor:integration:nonnegative-simple}，$c\mu(\{f\geqslant c\})\leqslant\int f\dd\mu$。

(2) 若 $f=0$ a.e.，由命题~\ref{prop:integration:nonnegative-basic}(3)，$\int f\dd\mu=\int0\dd\mu=0$。反之，设 $\int f\dd\mu=0$。由 (1)，$\mu(\{f\geqslant1/n\})\leqslant n\cdot0=0$。由于 $f\geqslant0$，$\{f\neq0\}=\{f>0\}=\bigcup_n\{f\geqslant1/n\}$ 是可数个零测集的并，为零测集。

(3) 对一切 $c>0$，$\mu(\{f=\infty\})\leqslant\mu(\{f\geqslant c\})\leqslant\frac1c\int f\dd\mu$，令 $c\to\infty$。
\end{proof}

\begin{example}[计数测度与级数]
取 $\Omega=\N$，$\cF=2^\N$，$\mu$ 为计数测度（它是 $\sigma$-有限的）。任意 $f:\N\to[0,\infty]$ 都可以写成 $f=\sum_if(i)\one_{\{i\}}$，由逐项积分与 $\int\one_{\{i\}}\dd\mu=1$，
\[
\int f\dd\mu=\sum_{i=1}^\infty f(i).
\]
所以非负项级数是积分的特例。对这个测度，单调收敛定理说的是：若 $0\leqslant a_{n,i}\uparrow a_i$（$n\to\infty$），则 $\sum_ia_{n,i}\uparrow\sum_ia_i$。逐项积分则说的是非负二重级数可以交换求和次序。这两个结论在数学分析中需要单独证明，这里成了一般定理的特例。
\end{example}

\topic{可积函数}

第四步，考虑一般的可测函数 $f:\Omega\to\R$。记 $f^+=\max\{f,0\}$，$f^-=\max\{-f,0\}$，则
\[
f=f^+-f^-,\qquad\abs f=f^++f^-,
\]
并且 $f^+,f^-$ 是非负可测函数（推论~\ref{cor:random-variables:arithmetic}）。

\begin{definition}[可积函数]\label{def:integration:integrable}
若 $\int\abs f\dd\mu<\infty$，称 $f$ \term{可积}（integrable），定义
\[
\int f\dd\mu=\int f^+\dd\mu-\int f^-\dd\mu .
\]
\end{definition}

由于 $f^\pm\leqslant\abs f$，两个积分都有限，右边有意义。当 $f$ 非负时，$f^-=0$，新的定义与第三步一致。我们顺便约定：若 $\int f^+\dd\mu$ 与 $\int f^-\dd\mu$ 中至少有一个有限，也用同一个公式定义 $\int f\dd\mu\in[-\infty,\infty]$。以后讨论期望时会用到这个约定。取值于 $\overline\R$ 的可积函数由命题~\ref{prop:integration:markov}(3) 几乎处处有限，把它在零测集 $\{\abs f=\infty\}$ 上改成 $0$ 不改变积分（命题~\ref{prop:integration:nonnegative-basic}(3)），所以下面只讨论实值函数。

\begin{lemma}\label{lem:integration:difference}
设 $f=f_1-f_2$，其中 $f_1,f_2\geqslant0$ 可积。则 $\int f\dd\mu=\int f_1\dd\mu-\int f_2\dd\mu$。
\end{lemma}
\begin{proof}
由 $f^+-f^-=f_1-f_2$ 得 $f_1+f^-=f_2+f^+$，四个函数都非负。由定理~\ref{thm:integration:linear-nonnegative}，
\[
\int f_1\dd\mu+\int f^-\dd\mu=\int f_2\dd\mu+\int f^+\dd\mu .
\]
各项都有限，移项即得。
\end{proof}

这个引理说，把 $f$ 写成两个非负可积函数之差，不管怎样写，算出的积分都一样。

\begin{theorem}\label{thm:integration:integrable}
设 $f,g$ 可积。则
\begin{enumerate}
\item 若 $f\geqslant0$ a.e.，则 $\int f\dd\mu\geqslant0$；
\item 对任意 $a\in\R$，$af$ 可积，且 $\int af\dd\mu=a\int f\dd\mu$；
\item $f+g$ 可积，且 $\int(f+g)\dd\mu=\int f\dd\mu+\int g\dd\mu$；
\item 若 $f\leqslant g$ a.e.，则 $\int f\dd\mu\leqslant\int g\dd\mu$；
\item 若 $f=g$ a.e.，则 $\int f\dd\mu=\int g\dd\mu$；
\item $\abs{\int f\dd\mu}\leqslant\int\abs f\dd\mu$。
\end{enumerate}
\end{theorem}
\begin{proof}
(1) $f\geqslant0$ a.e. 蕴含 $f^-=0$ a.e.，由命题~\ref{prop:integration:nonnegative-basic}(3)，$\int f^-\dd\mu=0$，于是 $\int f\dd\mu=\int f^+\dd\mu\geqslant0$。

(2) $\abs{af}=\abs a\abs f$，由命题~\ref{prop:integration:nonnegative-basic}(2) 可积。若 $a\geqslant0$，则 $(af)^+=af^+$，$(af)^-=af^-$，于是 $\int af\dd\mu=a\int f^+\dd\mu-a\int f^-\dd\mu=a\int f\dd\mu$。若 $a<0$，则 $(af)^+=\abs af^-$，$(af)^-=\abs af^+$，于是 $\int af\dd\mu=\abs a\int f^-\dd\mu-\abs a\int f^+\dd\mu=a\int f\dd\mu$。

(3) $\abs{f+g}\leqslant\abs f+\abs g$，由命题~\ref{prop:integration:nonnegative-basic}(1) 与定理~\ref{thm:integration:linear-nonnegative}，$\int\abs{f+g}\dd\mu\leqslant\int\abs f\dd\mu+\int\abs g\dd\mu<\infty$。又 $f+g=(f^++g^+)-(f^-+g^-)$，两个括号都是非负可积函数，由引理~\ref{lem:integration:difference} 与定理~\ref{thm:integration:linear-nonnegative}，
\[
\int(f+g)\dd\mu=\int f^+\dd\mu+\int g^+\dd\mu-\int f^-\dd\mu-\int g^-\dd\mu=\int f\dd\mu+\int g\dd\mu .
\]

(4)(5)(6) 由 (1)(2)(3) 与引理~\ref{lem:integration:consequences} 得到（可积函数的全体对加法、数乘、取绝对值封闭）。
\end{proof}

我们把常用的记号列在这里。
\begin{enumerate}
\item $(\Omega,\cF,\mu)=(\R,\cR,\lambda)$ 时，$\int f\dd\lambda$ 也写作 $\int f(x)\dd x$；对 $a<b$，$\int_{[a,b]}f\dd\lambda$ 写作 $\int_a^bf(x)\dd x$。下一节说明这与 Riemann 积分的记号是相容的。
\item $\mu=\mu_F$ 是 Stieltjes 测度函数 $F$ 对应的 Lebesgue--Stieltjes 测度时，$\int f\dd\mu_F$ 也写作 $\int f(x)\dd F(x)$，称为 \term{Lebesgue--Stieltjes 积分}。
\item $\Omega$ 可数，$\mu$ 是计数测度时，$\int f\dd\mu$ 就是 $\sum_{i\in\Omega}f(i)$（对非负 $f$ 见上面的例子；对可积的 $f$ 分成正、负部即可）。
\end{enumerate}

\topic{Lebesgue 积分与 Riemann 积分}

最后我们说明，对数学分析中遇到的函数，新的积分与 Riemann 积分给出同样的值。所以数学分析中计算积分的一切技巧，比如 Newton--Leibniz 公式、分部积分、换元，都可以继续用来计算 Lebesgue 积分。

\begin{theorem}\label{thm:integration:riemann}
设 $f:[a,b]\to\R$ 有界，Borel 可测，并且在 $[a,b]$ 上 Riemann 可积。则 $f\one_{[a,b]}$ 关于 $\lambda$ 可积，并且
\[
\int_{[a,b]}f\dd\lambda=\int_a^bf(x)\dd x\quad(\text{右边是 Riemann 积分}).
\]
特别地，这对 $[a,b]$ 上的连续函数成立。
\end{theorem}
\begin{proof}
在 $[a,b]$ 之外令 $f=0$，则 $f$ 有界，在测度有限的集合 $[a,b]$ 之外为零，所以它的积分由第二步定义。任取 $[a,b]$ 的分割 $P$：$a=x_0<x_1<\cdots<x_n=b$。令 $I_1=[x_0,x_1]$，$I_k=(x_{k-1},x_k]$（$2\leqslant k\leqslant n$），它们两两不交，并为 $[a,b]$。令 $m_k=\inf_{[x_{k-1},x_k]}f$，$M_k=\sup_{[x_{k-1},x_k]}f$，以及
\[
s_P=\sum_{k=1}^nm_k\one_{I_k},\qquad t_P=\sum_{k=1}^nM_k\one_{I_k}.
\]
它们是在 $[a,b]^c$ 上为零的简单函数，并且 $s_P\leqslant f\leqslant t_P$。由引理~\ref{lem:integration:representation} 与命题~\ref{prop:extension:lebesgue}(1)（$\lambda(I_k)=x_k-x_{k-1}$），
\[
\int s_P\dd\lambda=\sum_km_k(x_k-x_{k-1})=L(P),\qquad\int t_P\dd\lambda=\sum_kM_k(x_k-x_{k-1})=U(P),
\]
正是 Darboux 下和与上和。由定理~\ref{thm:integration:bounded} 中积分的上、下确界表示，$L(P)\leqslant\int f\dd\lambda\leqslant U(P)$ 对一切分割 $P$ 成立。$f$ Riemann 可积，意味着 $\sup_PL(P)=\inf_PU(P)=\int_a^bf(x)\dd x$，于是二者相等。
\end{proof}

\begin{remark}
定理中“Borel 可测”这个条件其实可以去掉，不过要把 Borel $\sigma$-代数换成 Lebesgue 可测集的 $\sigma$-代数：Riemann 可积函数都是 Lebesgue 可测的。这需要 Lebesgue 给出的 Riemann 可积性判别法：有界函数 $f$ 在 $[a,b]$ 上 Riemann 可积，当且仅当它的间断点集是 Lebesgue 零测集。这是本讲作业 B。反过来，Lebesgue 可积的函数远比 Riemann 可积的函数多，Dirichlet 函数 $\one_{\Q\cap[0,1]}$ 就是一个例子：它是简单函数，由推论~\ref{cor:integration:nonnegative-simple}，$\int\one_{\Q\cap[0,1]}\dd\lambda=\lambda(\Q\cap[0,1])=0$。
\end{remark}

对无界区间上的非负函数，单调收敛定理把 Lebesgue 积分与反常积分联系起来。设 $f\geqslant0$ 在 $\R$ 上连续。对 $n\geqslant1$，$f\one_{[-n,n]}\uparrow f$，由单调收敛定理与定理~\ref{thm:integration:riemann}，
\[
\int f\dd\lambda=\lim_{n\to\infty}\int_{-n}^nf(x)\dd x,
\]
右边是反常积分（两边都可以是 $\infty$）。对 $(-\infty,x]$ 上的积分同理。所以对第~\ref{lec:distributions} 讲中的密度 $f$（取为连续或分段连续的函数），$F(x)=\int_{-\infty}^xf(y)\dd y$ 中的积分按 Riemann 的反常积分理解与按 Lebesgue 积分 $\int_{(-\infty,x]}f\dd\lambda$ 理解是一样的。

对变号的函数要小心。$\int_0^\infty\frac{\sin x}x\dd x$ 作为反常积分收敛（等于 $\pi/2$），但 $\frac{\sin x}x$ 在 $(0,\infty)$ 上不是 Lebesgue 可积的，因为 $\int_0^\infty\frac{\abs{\sin x}}x\dd x=\infty$（习题 A9）。Lebesgue 积分要求绝对可积，它不承认靠正负相消得到的“条件收敛”。这与级数的情形一样：计数测度下的可积性就是绝对收敛。

\begin{workedexample}\label{we:integration:four-step}
设 $\mu,\nu$ 是 $(\Omega,\cF)$ 上的 $\sigma$-有限测度。用“四步法”证明：对非负可测函数 $f$，$\int f\dd(\mu+\nu)=\int f\dd\mu+\int f\dd\nu$；并由此计算 $\int_{[0,1]}x^2\dd(\lambda+\delta_1)(x)$。
\end{workedexample}
\begin{solution}
$\mu+\nu$ 是测度（\egref{ex:measures:combination}），并且是 $\sigma$-有限的（取 $E_n\uparrow\Omega$ 使 $\mu(E_n)<\infty$，$E_n'\uparrow\Omega$ 使 $\nu(E_n')<\infty$，则 $E_n\cap E_n'\uparrow\Omega$，两个测度都有限）。

\emph{第一步：示性函数。}由推论~\ref{cor:integration:nonnegative-simple}，$\int\one_A\dd(\mu+\nu)=(\mu+\nu)(A)=\mu(A)+\nu(A)=\int\one_A\dd\mu+\int\one_A\dd\nu$。

\emph{第二步：非负简单函数。}$\varphi=\sum_jy_j\one_{A_j}$，$y_j\geqslant0$。三个积分关于 $\varphi$ 都是线性的（定理~\ref{thm:integration:linear-nonnegative}），由第一步即得。

\emph{第三步：非负可测函数。}取非负简单函数 $\varphi_n\uparrow f$。由单调收敛定理（对三个测度各用一次），
\[
\int f\dd(\mu+\nu)=\lim_n\int\varphi_n\dd(\mu+\nu)=\lim_n\Bigl(\int\varphi_n\dd\mu+\int\varphi_n\dd\nu\Bigr)=\int f\dd\mu+\int f\dd\nu .
\]
（第四步，可积函数，拆成正部与负部即可，这里不需要。）

\emph{计算。}$\int_{[0,1]}x^2\dd\lambda=\frac13$（定理~\ref{thm:integration:riemann}），$\int_{[0,1]}x^2\dd\delta_1=1^2=1$（对 Dirac 测度，积分就是取值，见习题~\exref{ex:integration:A2}）。所以答案是 $\frac43$。如果把 $\lambda+\delta_1$ 限制在 $[0,1]$ 上再除以 $2$，就得到一个“一半均匀分布、一半集中在点 $1$”的概率分布，$\frac23$ 就是服从这个分布的随机变量 $X$ 的 $EX^2$。
\end{solution}

\begin{workedexample}\label{we:integration:series-telescoping}
用逐项积分计算 $\displaystyle\sum_{n=0}^\infty\frac1{(n+1)(n+2)}$，并计算 $\displaystyle\int_0^1\frac{x}{1-x}\log\frac1x\dd x$ 的值是否有限。
\end{workedexample}
\begin{solution}
\emph{第一问。}$\int_0^1x^n(1-x)\dd x=\frac1{n+1}-\frac1{n+2}=\frac1{(n+1)(n+2)}$（Riemann 积分，由定理~\ref{thm:integration:riemann} 也是 Lebesgue 积分）。函数 $f_n(x)=x^n(1-x)\one_{[0,1)}(x)$ 非负可测，并且对 $x\in[0,1)$，$\sum_nx^n(1-x)=(1-x)\cdot\frac1{1-x}=1$。由逐项积分（推论~\ref{cor:integration:series}），
\[
\sum_{n=0}^\infty\frac1{(n+1)(n+2)}=\sum_{n=0}^\infty\int f_n\dd\lambda=\int\sum_{n=0}^\infty f_n\dd\lambda=\int\one_{[0,1)}\dd\lambda=1 .
\]
这里不需要任何一致收敛：$\sum_nx^n(1-x)$ 在 $[0,1)$ 上并不一致收敛（每个部分和在 $x\to1$ 时趋于 $0$，而和是 $1$）。

\emph{第二问。}对 $x\in(0,1)$，$\frac x{1-x}=\sum_{n\geqslant1}x^n$，各项非负，所以
\[
\int_0^1\frac x{1-x}\log\frac1x\dd x=\sum_{n=1}^\infty\int_0^1x^n\log\frac1x\dd x=\sum_{n=1}^\infty\frac1{(n+1)^2}<\infty .
\]
其中 $\int_0^1x^n\log\frac1x\dd x=\frac1{(n+1)^2}$ 由分部积分得到（$[\frac{x^{n+1}}{n+1}\log\frac1x]_0^1+\int_0^1\frac{x^n}{n+1}\dd x$；在 $0$ 附近的反常积分由单调收敛定理处理，与上面关于反常积分的说明相同）。所以积分有限，值为 $\sum_{k\geqslant2}\frac1{k^2}=\frac{\pi^2}6-1$。
\end{solution}

\topic{Fatou 引理}

单调收敛定理只管单调的函数列。可是在概率论中，我们遇到的函数列多半不单调，比如样本均值 $(X_1+\cdots+X_n)/n$，它随 $n$ 忽大忽小。所以我们需要能处理一般函数列的极限定理。这一节和下一节给出两个：Fatou 引理与控制收敛定理。它们都是从单调收敛定理推出来的。

我们先看几个例子，弄清楚“积分的极限”与“极限的积分”可以差多少。

\begin{example}[质量的三种逃逸方式]\label{ex:integration:escape}
下面三个例子中，$f_n\geqslant0$，$f_n\to0$ 处处成立，可是 $\int f_n\dd\lambda=1$ 对一切 $n$ 成立，所以
\[
\int\lim_nf_n\dd\lambda=0<1=\lim_n\int f_n\dd\lambda .
\]
\begin{enumerate}
\item \emph{向上逃逸。}在 $(0,1)$ 上取 $f_n=n\one_{(0,1/n)}$。图像是一个底为 $1/n$、高为 $n$ 的长方形，面积为 $1$。对固定的 $x\in(0,1)$，当 $n>1/x$ 时 $x\notin(0,1/n)$，所以 $f_n(x)=0$。长方形越来越瘦，越来越高，最后“缩”到 $x=0$ 这一点上，而这一点不在 $(0,1)$ 中。
\item \emph{向右逃逸。}在 $\R$ 上取 $f_n=\one_{[n,n+1]}$。长方形的形状不变，但是一直往右走。对固定的 $x$，当 $n>x$ 时 $f_n(x)=0$。
\item \emph{摊平。}在 $\R$ 上取 $f_n=\frac1n\one_{[0,n]}$。长方形越来越矮，越来越宽。这里 $f_n\to0$ 甚至是一致的：$\sup_x\abs{f_n(x)}=1/n$。
\end{enumerate}
这三个例子说明，一般的函数列，极限与积分不能交换；即使一致收敛也不行（第三个例子，在测度无限的空间上）。它们还说明了一个现象：在极限过程中，积分只会“丢失”，而不会“凭空产生”。三个例子中丢失的都是 $1$，极限函数的积分比积分的极限小。Fatou 引理说，对非负函数这总是对的。
\end{example}

\begin{theorem}[Fatou 引理]\label{thm:integration:fatou}
设 $f_n:\Omega\to[0,\infty]$ 可测。则
\[
\int\liminf_{n\to\infty}f_n\dd\mu\leqslant\liminf_{n\to\infty}\int f_n\dd\mu .
\]
\end{theorem}

\begin{remark}[证明思路]
$\liminf_nf_n$ 按定义是 $\lim_n\inf_{m\geqslant n}f_m$，而 $g_n=\inf_{m\geqslant n}f_m$ 关于 $n$ 单调递增。所以左边可以用单调收敛定理算出来。另一方面，$g_n$ 不超过每个 $f_m$（$m\geqslant n$），它的积分就不超过每个 $\int f_m$。两件事合起来就是结论。
\end{remark}

\begin{proof}
令 $g_n=\inf_{m\geqslant n}f_m$。

\emph{第一步：$g_n$ 非负可测，并且单调递增地收敛到 $\liminf_nf_n$。}由定理~\ref{thm:random-variables:limits}，$g_n$ 是可测的，取值于 $[0,\infty]$。当 $n$ 增大时，取下确界的范围 $\{m:m\geqslant n\}$ 变小，所以 $g_n\leqslant g_{n+1}$。按照下极限的定义，$\liminf_nf_n=\lim_ng_n=\sup_ng_n$。

\emph{第二步：比较积分。}对每个 $m\geqslant n$，$g_n\leqslant f_m$，由命题~\ref{prop:integration:nonnegative-basic}(1)，$\int g_n\dd\mu\leqslant\int f_m\dd\mu$。对 $m\geqslant n$ 取下确界，
\[
\int g_n\dd\mu\leqslant\inf_{m\geqslant n}\int f_m\dd\mu .
\]

\emph{第三步：取极限。}令 $n\to\infty$。由第一步与单调收敛定理（定理~\ref{thm:integration:mct}），左边收敛到 $\int\liminf_nf_n\dd\mu$；右边是数列 $\int f_m\dd\mu$ 的“后段下确界”，单调递增地收敛到 $\liminf_m\int f_m\dd\mu$。不等号在极限中保持，结论得证。
\end{proof}

我们来看 Fatou 引理的条件起什么作用，以及它和前面学过的东西的联系。

\begin{example}[非负性不能去掉]
在 $(0,1)$ 上取 $f_n=-n\one_{(0,1/n)}$。则 $f_n\to0$ 处处成立，$\int\liminf_nf_n\dd\lambda=0$，而 $\int f_n\dd\lambda=-1$，于是 $\liminf_n\int f_n\dd\lambda=-1<0$。Fatou 引理的不等号反过来了。这就是\egref{ex:integration:escape} 中第一个例子的相反数：对非负函数积分只会丢失，对非正函数积分只会“凭空增加”。
\end{example}

\begin{example}[不等号可以是严格的，并且右边不能换成上极限]
设 $A\in\cF$，令 $f_n=\one_A$（$n$ 为偶数），$f_n=\one_{A^c}$（$n$ 为奇数）。每个点要么属于 $A$ 要么属于 $A^c$，所以 $\liminf_nf_n=0$。而 $\int f_n\dd\mu$ 在 $\mu(A)$ 与 $\mu(A^c)$ 之间来回跳。在 $([0,1],\lambda)$ 上取 $A=[0,1/3]$，则
\[
\int\liminf_nf_n\dd\lambda=0<\frac13=\liminf_n\int f_n\dd\lambda<\frac23=\limsup_n\int f_n\dd\lambda .
\]
\end{example}

\begin{example}[集合的 Fatou 引理]
取 $f_n=\one_{A_n}$。由于 $\one_{A_n}(\omega)$ 只取 $0$ 和 $1$，$\liminf_n\one_{A_n}(\omega)=1$ 当且仅当 $\omega$ 从某项起一直属于 $A_n$，也就是 $\omega\in\liminf_nA_n$（定义~\ref{def:set-systems:limsup}）。所以 $\liminf_n\one_{A_n}=\one_{\liminf_nA_n}$，Fatou 引理给出
\[
\mu\bigl(\liminf_nA_n\bigr)\leqslant\liminf_n\mu(A_n),
\]
这就是命题~\ref{prop:measures:fatou-sets}(1)。所以 Fatou 引理是那个关于集合的命题在函数上的推广，这也是它名字的来历。
\end{example}

Fatou 引理最常见的用法是：已知 $f_n\to f$ a.e.，并且 $\int f_n\dd\mu\leqslant C$ 对一切 $n$ 成立，那么 $\int f\dd\mu\leqslant C$。也就是说，\emph{积分的有界性在几乎处处收敛的极限下保持}。我们把它写成一个推论。

\begin{corollary}\label{cor:integration:fatou-bound}
设 $f_n,f:\Omega\to[0,\infty]$ 可测，$f_n\to f$ a.e.，并且 $\sup_n\int f_n\dd\mu\leqslant C$。则 $\int f\dd\mu\leqslant C$。
\end{corollary}
\begin{proof}
设 $N$ 是零测集，在 $N^c$ 上 $f_n\to f$。在 $N^c$ 上 $\liminf_nf_n=f$，所以 $\liminf_nf_n=f$ a.e.。由命题~\ref{prop:integration:nonnegative-basic}(3) 与 Fatou 引理，
\[
\int f\dd\mu=\int\liminf_nf_n\dd\mu\leqslant\liminf_n\int f_n\dd\mu\leqslant C .\qedhere
\]
\end{proof}

\topic{控制收敛定理}

Fatou 引理只给出一个不等式。要得到等式，就得排除\egref{ex:integration:escape} 中那样的“逃逸”。在那三个例子中，没有一个可积函数能把全部 $f_n$ 都压在下面：第一个例子中 $\sup_nf_n(x)$ 大约是 $1/x$，在 $(0,1)$ 上不可积；第二个例子中 $\sup_nf_n=\one_{[1,\infty)}$，第三个例子中 $\sup_nf_n(x)$ 大约是 $\min\{1,1/x\}$，在 $\R$ 上都不可积。下面的定理说，只要存在这样一个可积的“控制函数”，质量就逃不掉，极限与积分可以交换。

\begin{theorem}[Lebesgue 控制收敛定理]\label{thm:integration:dct}
设 $f_n,f:\Omega\to\R$ 可测，$f_n\to f$ a.e.。若存在可积函数 $g$，使对每个 $n$，$\abs{f_n}\leqslant g$ a.e.，则 $f$ 可积，并且
\[
\lim_{n\to\infty}\int\abs{f_n-f}\dd\mu=0,\qquad\text{从而}\qquad\lim_{n\to\infty}\int f_n\dd\mu=\int f\dd\mu .
\]
\end{theorem}

\begin{remark}[证明思路]
要证明 $\int\abs{f_n-f}\to0$。$\abs{f_n-f}$ 不超过 $2g$，于是 $2g-\abs{f_n-f}$ 是非负的，并且收敛到 $2g$。对它用 Fatou 引理，得到 $\int2g\leqslant\int2g-\limsup\int\abs{f_n-f}$。因为 $\int2g$ 是有限数，可以从两边消去，就得到 $\limsup\int\abs{f_n-f}\leqslant0$。这里 $g$ 可积是必不可少的：如果 $\int g=\infty$，就无法从两边消去 $\infty$。
\end{remark}

\begin{proof}
\emph{第一步：把“几乎处处”化为“处处”。}对每个 $n$，有零测集 $N_n$，在 $N_n^c$ 上 $\abs{f_n}\leqslant g$；又有零测集 $N_0$，在 $N_0^c$ 上 $f_n\to f$。可积函数 $g$ 几乎处处有限（命题~\ref{prop:integration:markov}(3)），设 $\{g=\infty\}=N'$。令 $N=N_0\cup N'\cup\bigcup_{n\geqslant1}N_n$，它是可数个零测集的并，仍是零测集。把 $f_n,f,g$ 都换成 $f_n\one_{N^c}$，$f\one_{N^c}$，$g\one_{N^c}$。它们仍然可测；由命题~\ref{prop:integration:nonnegative-basic}(3) 与定理~\ref{thm:integration:integrable}(5)，这样修改不改变任何一个积分。所以下面不妨设：处处有 $f_n\to f$，$\abs{f_n}\leqslant g<\infty$。

\emph{第二步：$f$ 可积。}在 $\abs{f_n}\leqslant g$ 中令 $n\to\infty$，得 $\abs f\leqslant g$。于是 $\int\abs f\dd\mu\leqslant\int g\dd\mu<\infty$。

\emph{第三步：用 Fatou 引理。}令 $h_n=2g-\abs{f_n-f}$。由 $\abs{f_n-f}\leqslant\abs{f_n}+\abs f\leqslant2g$ 知 $h_n\geqslant0$；它可测；并且由于 $f_n\to f$，$h_n\to2g$。对 $h_n$ 用 Fatou 引理（定理~\ref{thm:integration:fatou}）：
\[
\int2g\dd\mu=\int\liminf_nh_n\dd\mu\leqslant\liminf_n\int h_n\dd\mu .
\]
$\abs{f_n-f}$ 与 $2g$ 都可积，由线性性（定理~\ref{thm:integration:integrable}(2)(3)），$\int h_n\dd\mu=\int2g\dd\mu-\int\abs{f_n-f}\dd\mu$。于是
\[
\int2g\dd\mu\leqslant\liminf_n\Bigl(\int2g\dd\mu-\int\abs{f_n-f}\dd\mu\Bigr)=\int2g\dd\mu-\limsup_n\int\abs{f_n-f}\dd\mu .
\]

\emph{第四步：消去。}$\int2g\dd\mu$ 是有限数，从两边消去，得 $\limsup_n\int\abs{f_n-f}\dd\mu\leqslant0$。而 $\int\abs{f_n-f}\dd\mu\geqslant0$，所以它的极限存在且等于 $0$。最后，由定理~\ref{thm:integration:integrable}(3)(6)，
\[
\Bigl\lvert\int f_n\dd\mu-\int f\dd\mu\Bigr\rvert=\Bigl\lvert\int(f_n-f)\dd\mu\Bigr\rvert\leqslant\int\abs{f_n-f}\dd\mu\to0 .\qedhere
\]
\end{proof}

\begin{remark}[控制函数的含义]
条件“$\abs{f_n}\leqslant g$ 对一切 $n$ 成立，$g$ 可积”等价于说 $\sup_n\abs{f_n}$ 被一个可积函数控制，也就是 $\int\sup_n\abs{f_n}\dd\mu<\infty$（若这个积分有限，取 $g=\sup_n\abs{f_n}$ 即可）。直观上，它规定了一个所有 $f_n$ 都不许越出的“围栏”，而这个围栏本身的面积有限。\egref{ex:integration:escape} 的三个例子都越出了任何有限面积的围栏。控制收敛定理的条件是充分的，但不是必要的；以后讲一致可积性时，我们会看到在概率空间上什么样的条件才是“恰好”的。
\end{remark}

在概率空间或者一般的有限测度空间上，常数是可积的，于是有下面这个常用的特例。

\begin{corollary}[有界收敛定理]\label{cor:integration:bct}
设 $\mu(\Omega)<\infty$，$f_n\to f$ a.e.，并且存在常数 $M$ 使 $\abs{f_n}\leqslant M$ a.e. 对一切 $n$ 成立。则 $\int\abs{f_n-f}\dd\mu\to0$，$\int f_n\dd\mu\to\int f\dd\mu$。
\end{corollary}
\begin{proof}
取 $g=M$（常函数），$\int g\dd\mu=M\mu(\Omega)<\infty$，由控制收敛定理即得。
\end{proof}

\begin{corollary}[可积函数项级数的逐项积分]\label{cor:integration:series-integrable}
设 $f_n$ 可测，$\sum_{n=1}^\infty\int\abs{f_n}\dd\mu<\infty$。则级数 $\sum_nf_n$ 几乎处处绝对收敛，其和（在不收敛的零测集上定义为 $0$）可积，并且
\[
\int\sum_{n=1}^\infty f_n\dd\mu=\sum_{n=1}^\infty\int f_n\dd\mu .
\]
\end{corollary}
\begin{proof}
令 $g=\sum_n\abs{f_n}$，取值于 $[0,\infty]$。由逐项积分（推论~\ref{cor:integration:series}），$\int g\dd\mu=\sum_n\int\abs{f_n}\dd\mu<\infty$，所以 $g$ 可积，从而 $g<\infty$ a.e.（命题~\ref{prop:integration:markov}(3)）。在 $\{g<\infty\}$ 上，级数 $\sum_nf_n$ 绝对收敛。令 $N=\{g=\infty\}$，部分和 $s_N=\sum_{n\leqslant N}f_n\one_{N^c}$。则 $s_N\to s=\sum_nf_n\one_{N^c}$ 处处成立，并且 $\abs{s_N}\leqslant g$。由控制收敛定理，$s$ 可积，并且 $\int s\dd\mu=\lim_N\int s_N\dd\mu=\lim_N\sum_{n\leqslant N}\int f_n\dd\mu$（有限和的积分可以拆开，这是线性性；在零测集 $N$ 上的修改不改变积分）。
\end{proof}

下面是控制收敛定理的几个典型应用。做这类题目的步骤总是一样的，同学们应该熟记：（1）写出被积函数列 $f_n$，求出逐点极限 $f$；（2）找一个与 $n$ 无关的可积函数 $g$，使 $\abs{f_n}\leqslant g$；（3）由控制收敛定理交换极限与积分。其中第（2）步是需要动脑筋的地方。

\begin{workedexample}\label{we:integration:exp-limit}
证明 $\displaystyle\lim_{n\to\infty}\int_0^n\Bigl(1-\frac xn\Bigr)^n\dd x=1$。
\end{workedexample}
\begin{solution}
令 $f_n(x)=(1-x/n)^n\one_{[0,n]}(x)$，$x\in[0,\infty)$。这样 $\int_0^n(1-x/n)^n\dd x=\int_{[0,\infty)}f_n\dd\lambda$，积分区域就不再随 $n$ 变化了。（每个 $f_n$ 在 $[0,n]$ 上连续，在 $(n,\infty)$ 上为零，是 Borel 函数；它在 $[0,n]$ 上的 Lebesgue 积分等于 Riemann 积分，由定理~\ref{thm:integration:riemann}。）

\emph{逐点极限。}固定 $x\geqslant0$。当 $n>x$ 时 $x\in[0,n]$，$f_n(x)=(1-x/n)^n$。由数学分析中的重要极限，$(1-x/n)^n\to e^{-x}$。所以 $f_n(x)\to e^{-x}$。

\emph{控制函数。}对一切实数 $y$，有 $1-y\leqslant e^{-y}$（函数 $e^{-y}-1+y$ 在 $y=0$ 处取最小值 $0$）。当 $0\leqslant x\leqslant n$ 时，$0\leqslant1-x/n\leqslant e^{-x/n}$，两边 $n$ 次方得 $0\leqslant f_n(x)\leqslant e^{-x}$；当 $x>n$ 时 $f_n(x)=0\leqslant e^{-x}$。所以 $\abs{f_n}\leqslant g$，$g(x)=e^{-x}$。$g$ 在 $[0,\infty)$ 上可积：由单调收敛定理与定理~\ref{thm:integration:riemann}，$\int_{[0,\infty)}e^{-x}\dd x=\lim_{N}\int_0^Ne^{-x}\dd x=\lim_N(1-e^{-N})=1$。

\emph{交换极限与积分。}由控制收敛定理，$\int f_n\dd\lambda\to\int_{[0,\infty)}e^{-x}\dd x=1$。
\end{solution}

\begin{workedexample}\label{we:integration:dct-fails}
在 $(0,1)$ 上取 $f_n(x)=n^2x(1-x)^n$。求 $\lim_nf_n(x)$ 与 $\lim_n\int_0^1f_n(x)\dd x$。能否找到控制函数？
\end{workedexample}
\begin{solution}
\emph{逐点极限。}固定 $x\in(0,1)$，令 $q=1-x\in(0,1)$。则 $f_n(x)=xn^2q^n\to0$，因为 $n^2q^n\to0$（指数衰减快于多项式增长，例如由 $\sum_nn^2q^n$ 收敛的比值判别法）。所以 $f_n\to0$ 处处成立。

\emph{积分。}用分部积分或 Beta 积分：$\int_0^1x(1-x)^n\dd x$，令 $u=1-x$，等于 $\int_0^1(1-u)u^n\dd u=\frac1{n+1}-\frac1{n+2}=\frac1{(n+1)(n+2)}$。于是
\[
\int_0^1f_n(x)\dd x=\frac{n^2}{(n+1)(n+2)}\to1\neq0 .
\]
极限与积分不能交换。

\emph{控制函数不存在。}如果存在可积的 $g$ 使 $\abs{f_n}\leqslant g$ a.e. 对一切 $n$ 成立，控制收敛定理就会给出 $\int f_n\to0$，与上面的计算矛盾。所以不存在这样的 $g$。我们也可以直接看出 $\sup_nf_n$ 不可积：当 $n\geqslant2$，$x\in[\frac1{n+1},\frac1n]$ 时，$(1-x)^n\geqslant(1-\frac1n)^n\geqslant\frac14$（数列 $(1-\frac1n)^n$ 单调递增，$n=2$ 时等于 $\frac14$），所以
\[
f_n(x)\geqslant n^2\cdot\frac1{n+1}\cdot\frac14\geqslant\frac{n+1}{16}\geqslant\frac1{16x},
\]
其中用到 $n^2\geqslant(n+1)^2/4$ 与 $n+1\geqslant1/x$。这些区间盖住了 $(0,\frac12]$，所以 $\sup_nf_n(x)\geqslant\frac1{16x}$ 在 $(0,\frac12]$ 上成立，而 $\int_0^{1/2}\frac{\dd x}x=\infty$。这正是\egref{ex:integration:escape} 中“向上逃逸”的情形：$f_n$ 的质量挤到了 $x=0$ 附近。
\end{solution}

\begin{workedexample}\label{we:integration:zeta2}
证明 $\displaystyle\int_0^\infty\frac x{e^x-1}\dd x=\sum_{n=1}^\infty\frac1{n^2}$。
\end{workedexample}
\begin{solution}
思路是把被积函数展开成级数，再逐项积分。对 $x>0$，$0<e^{-x}<1$，由几何级数，
\[
\frac x{e^x-1}=\frac{xe^{-x}}{1-e^{-x}}=\sum_{n=1}^\infty xe^{-nx}.
\]
每一项 $xe^{-nx}$ 非负、连续，所以可以用非负函数的逐项积分（推论~\ref{cor:integration:series}），不需要任何控制：
\[
\int_{(0,\infty)}\frac x{e^x-1}\dd x=\sum_{n=1}^\infty\int_{(0,\infty)}xe^{-nx}\dd x .
\]
对每个 $n$，由单调收敛定理与定理~\ref{thm:integration:riemann}，$\int_{(0,\infty)}xe^{-nx}\dd x=\lim_N\int_0^Nxe^{-nx}\dd x$。分部积分，
\[
\int_0^Nxe^{-nx}\dd x=\Bigl[-\frac{x}{n}e^{-nx}\Bigr]_0^N+\frac1n\int_0^Ne^{-nx}\dd x=-\frac Nne^{-nN}+\frac1{n^2}(1-e^{-nN})\to\frac1{n^2}.
\]
所以原积分等于 $\sum_n1/n^2$。（以后在特征函数一讲会顺便知道这个和是 $\pi^2/6$。）

注意：这个例子里我们用的是非负项的逐项积分，它不需要控制函数，只需要每一项非负。这比数学分析中逐项积分需要一致收敛方便得多：这里的级数在 $x\to0$ 附近并不一致收敛。
\end{solution}

\topic{含参变量的积分}

控制收敛定理的一个重要应用是含参变量积分的连续性与求导。在数学分析中，这些结论需要一致收敛之类的条件；现在只需要一个控制函数。以后在特征函数一讲，我们会用它来对 $\varphi(t)=Ee^{itX}$ 求导。

\begin{theorem}[含参变量积分的连续性]\label{thm:integration:parameter-continuity}
设 $I\subseteq\R$ 是区间，$f:\Omega\times I\to\R$ 满足：
\begin{enumerate}
\item 对每个 $t\in I$，$x\mapsto f(x,t)$ 可测；
\item 对每个 $x\in\Omega$，$t\mapsto f(x,t)$ 在 $t_0\in I$ 处连续；
\item 存在可积函数 $g$，使对一切 $t\in I$，$\abs{f(x,t)}\leqslant g(x)$ 对一切 $x$ 成立。
\end{enumerate}
则 $F(t)=\int f(x,t)\,\mu(\dd x)$ 在 $t_0$ 处连续。
\end{theorem}
\begin{proof}
由 (3)，每个 $f(\cdot,t)$ 可积，$F(t)$ 有定义。实变函数在一点连续，等价于对任意 $t_n\to t_0$（$t_n\in I$）有 $F(t_n)\to F(t_0)$（数学分析中的 Heine 定理）。任取这样一列 $t_n$，令 $f_n(x)=f(x,t_n)$。由 (2)，$f_n(x)\to f(x,t_0)$ 对每个 $x$ 成立；由 (3)，$\abs{f_n}\leqslant g$。由控制收敛定理，$F(t_n)=\int f_n\dd\mu\to\int f(x,t_0)\,\mu(\dd x)=F(t_0)$。
\end{proof}

\begin{theorem}[积分号下求导]\label{thm:integration:differentiation}
设 $I\subseteq\R$ 是开区间，$f:\Omega\times I\to\R$ 满足：
\begin{enumerate}
\item 对每个 $t\in I$，$x\mapsto f(x,t)$ 可积；
\item 对每个 $x\in\Omega$，$t\mapsto f(x,t)$ 在 $I$ 上可导，导数记为 $\partial_tf(x,t)$；
\item 存在可积函数 $g$，使 $\abs{\partial_tf(x,t)}\leqslant g(x)$ 对一切 $x\in\Omega$，$t\in I$ 成立。
\end{enumerate}
则 $F(t)=\int f(x,t)\,\mu(\dd x)$ 在 $I$ 上可导，并且 $F'(t)=\int\partial_tf(x,t)\,\mu(\dd x)$。
\end{theorem}
\begin{proof}
固定 $t\in I$，任取 $h_n\to0$，$h_n\neq0$，$t+h_n\in I$。令
\[
q_n(x)=\frac{f(x,t+h_n)-f(x,t)}{h_n}.
\]
由 (1)，$q_n$ 可测且可积，并且 $\frac{F(t+h_n)-F(t)}{h_n}=\int q_n\dd\mu$。由 (2)，$q_n(x)\to\partial_tf(x,t)$ 对每个 $x$ 成立；特别地 $x\mapsto\partial_tf(x,t)$ 作为可测函数列的极限是可测的（定理~\ref{thm:random-variables:limits}）。由微分中值定理，对每个 $x$，存在 $\theta$ 介于 $t$ 与 $t+h_n$ 之间（$\theta$ 依赖于 $x$ 与 $n$），使 $q_n(x)=\partial_tf(x,\theta)$，于是由 (3)，$\abs{q_n(x)}\leqslant g(x)$。由控制收敛定理，
\[
\frac{F(t+h_n)-F(t)}{h_n}=\int q_n\dd\mu\to\int\partial_tf(x,t)\,\mu(\dd x).
\]
由于 $h_n$ 是任意趋于零的数列，由 Heine 定理，$F'(t)$ 存在并等于右边。
\end{proof}

\begin{remark}
证明中我们并不需要 $\theta$ 作为 $x$ 的函数可测，因为我们只是用它来逐点估计 $\abs{q_n(x)}$；$q_n$ 本身的可测性是由它的定义式得到的。这是这类证明中容易出错的地方。另外，条件 (3) 往往不能在整个 $I$ 上成立，这时只需对每个 $t_0\in I$，在 $t_0$ 的某个小开区间上验证，就得到 $F$ 在 $t_0$ 处可导。
\end{remark}

\begin{workedexample}\label{we:integration:factorial}
设 $n\geqslant0$ 是整数。证明 $\displaystyle\int_0^\infty x^ne^{-x}\dd x=n!$。
\end{workedexample}
\begin{solution}
思路：先算一个依赖参数 $t$ 的简单积分，再对 $t$ 反复求导，每求一次导数，被积函数就多出一个因子 $-x$。

对 $t>0$，令 $F(t)=\int_{(0,\infty)}e^{-tx}\dd x$。和\wegref{we:integration:exp-limit} 中一样，由单调收敛定理与 Riemann 积分，$F(t)=\lim_N\frac{1-e^{-tN}}t=\frac1t$。

我们用归纳法证明：对每个 $k\geqslant0$ 与 $t>0$，$x^ke^{-tx}$ 在 $(0,\infty)$ 上可积，并且
\[
\int_{(0,\infty)}x^ke^{-tx}\dd x=\frac{k!}{t^{k+1}} .
\]
$k=0$ 是上面的计算。设对 $k$ 成立，考虑 $f(x,t)=x^ke^{-tx}$。固定 $t_0>0$，在开区间 $I=(t_0/2,\infty)$ 上验证定理~\ref{thm:integration:differentiation} 的条件：(1) 由归纳假设，每个 $f(\cdot,t)$ 可积；(2) $\partial_tf(x,t)=-x^{k+1}e^{-tx}$；(3) 当 $t\in I$ 时，$\abs{\partial_tf(x,t)}\leqslant x^{k+1}e^{-t_0x/2}=:g(x)$，而 $g$ 可积，这是因为 $x^{k+1}e^{-t_0x/4}$ 在 $(0,\infty)$ 上有界（它是连续的，并且在 $x\to\infty$ 时趋于 $0$），设上界为 $C$，则 $g(x)\leqslant Ce^{-t_0x/4}$，而后者的积分是 $4C/t_0<\infty$。于是在 $t_0$ 处，
\[
-\int_{(0,\infty)}x^{k+1}e^{-t_0x}\dd x=\frac{\dd}{\dd t}\Bigl(\frac{k!}{t^{k+1}}\Bigr)\Big|_{t=t_0}=-\frac{(k+1)!}{t_0^{k+2}} ,
\]
并且左边的被积函数 $\leqslant g$ 可积。这就是 $k+1$ 时的结论。最后取 $t=1$。
\end{solution}

\topic{本讲小结}

这一讲建立了 Lebesgue 积分，我们把逻辑线索再理一遍。
\begin{enumerate}
\item 积分分四步定义：简单函数（直接按“值乘测度”求和），有界函数（用简单函数从上下两边夹，定理~\ref{thm:integration:bounded}），非负函数（用下方的有界函数取上确界，定义~\ref{def:integration:nonnegative}），可积函数（$\int f=\int f^+-\int f^-$）。
\item 整个理论的支柱是 Levi 单调收敛定理（定理~\ref{thm:integration:mct}）。它的证明用到测度的可数可加性，这是 Lebesgue 积分比 Riemann 积分好用的根本原因。
\item 由单调收敛定理推出：积分的线性性，非负函数逐项积分，不定积分是测度，Markov 不等式；再推出 Fatou 引理（对非负函数，积分在极限中只会丢失），由 Fatou 引理推出控制收敛定理（有可积的控制函数时，极限与积分可以交换）。
\item 证明关于积分的等式时，标准的套路是“四步法”：示性函数，简单函数，非负函数（单调收敛定理），可积函数（拆成正部与负部）。
\item 对数学分析中的函数，Lebesgue 积分与 Riemann 积分相同（定理~\ref{thm:integration:riemann}），所以具体的积分仍然用 Newton--Leibniz 公式计算。
\end{enumerate}
下一讲之后，我们把这些定理翻译成期望的语言：$EX=\int X\dd P$，单调收敛定理、Fatou 引理和控制收敛定理将是我们处理随机变量极限的基本工具。

% =====================================================================
\homework{Riemann 可积性的 Lebesgue 判别法，分层积分公式，Riemann--Lebesgue 引理与连续函数逼近}[Riemann 可积性的判别，分层积分公式，Riemann–Lebesgue 引理]
\label{hw:integration}

\exercisepart{基本习题}

\begin{exercises}{A}{课堂内容的补充}
\item\label{ex:integration:A1} 设 $f\geqslant0$ 可测，$E_{n,m}=\{m2^{-n}\leqslant f<(m+1)2^{-n}\}$。证明 $n\to\infty$ 时
\[
\sum_{m=1}^\infty\frac m{2^n}\mu(E_{n,m})\ \uparrow\ \int_{\{f<\infty\}}f\dd\mu ;
\]
特别地，若 $\mu(\{f=\infty\})=0$，则左边单调递增地收敛到 $\int f\dd\mu$。（提示：由逐项积分，左边是一个非负函数的积分，这个函数是什么？它关于 $n$ 单调吗？）
\begin{solution}
思路：把和式认出来是一个简单函数（可数个值的函数）的积分，然后用单调收敛定理。

令 $g_n=\sum_{m\geqslant1}\frac m{2^n}\one_{E_{n,m}}$。$E_{n,m}$（$m\geqslant0$）两两不交，并为 $\{f<\infty\}$，所以在 $\{f<\infty\}$ 上 $g_n=2^{-n}\lfloor2^nf\rfloor$，在 $\{f=\infty\}$ 上 $g_n=0$。由逐项积分（推论~\ref{cor:integration:series}）与 $\int\one_A\dd\mu=\mu(A)$，
\[
\int g_n\dd\mu=\sum_{m\geqslant1}\frac m{2^n}\mu(E_{n,m}).
\]
\emph{单调。}对 $t\geqslant0$，$2\lfloor t\rfloor\leqslant\lfloor2t\rfloor$，所以 $2^{-n}\lfloor2^nf\rfloor\leqslant2^{-n-1}\lfloor2^{n+1}f\rfloor$，$g_n\leqslant g_{n+1}$。

\emph{极限。}在 $\{f<\infty\}$ 上，$f-2^{-n}<g_n\leqslant f$，所以 $g_n\uparrow f\one_{\{f<\infty\}}$。由单调收敛定理，$\int g_n\dd\mu\uparrow\int f\one_{\{f<\infty\}}\dd\mu$。若 $\mu(\{f=\infty\})=0$，则 $f\one_{\{f<\infty\}}=f$ a.e.，两者积分相同（命题~\ref{prop:integration:nonnegative-basic}(3)）。

注意 $\mu(\{f=\infty\})>0$ 时 $\int f\dd\mu=\infty$，而左边可以是有限的（比如 $f=\infty\cdot\one_A$，左边恒为 $0$），所以这个限制是必要的。
\end{solution}

\item\label{ex:integration:A2} Dirac 测度与离散测度上的积分。
\begin{questions}
\item 设 $x\in\Omega$，$\delta_x$ 是 Dirac 测度。证明对非负可测函数 $f$，$\int f\dd\delta_x=f(x)$。
\begin{solution}
用四步法的前三步。

\emph{示性函数。}$\int\one_A\dd\delta_x=\delta_x(A)=\one_A(x)$。

\emph{非负简单函数。}$\varphi=\sum_jy_j\one_{A_j}$，由线性性，$\int\varphi\dd\delta_x=\sum_jy_j\one_{A_j}(x)=\varphi(x)$。

\emph{非负可测函数。}取非负简单函数 $\varphi_n\uparrow f$，由单调收敛定理，$\int f\dd\delta_x=\lim_n\int\varphi_n\dd\delta_x=\lim_n\varphi_n(x)=f(x)$。
\end{solution}
\item 设 $\mu=\sum_ip_i\delta_{x_i}$（$p_i\geqslant0$），证明 $\int f\dd\mu=\sum_ip_if(x_i)$。
\begin{solution}
本讲的积分只对 $\sigma$-有限测度定义，所以这里假定 $\mu$ 是 $\sigma$-有限的（例如 $\sum_ip_i<\infty$ 时 $\mu$ 是有限测度）。

示性函数：$\int\one_A\dd\mu=\mu(A)=\sum_ip_i\one_A(x_i)$。简单函数：由线性性与非负级数可以逐项相加。非负函数：$\varphi_n\uparrow f$，左边由单调收敛定理趋于 $\int f\dd\mu$；右边 $\sum_ip_i\varphi_n(x_i)\uparrow\sum_ip_if(x_i)$，这是计数测度上的单调收敛定理（\egref{ex:measures:counting} 之后关于级数的例子），或者直接由非负级数的性质得到。
\end{solution}
\end{questions}

\item\label{ex:integration:A3} 计数测度 $(\N,2^\N,\#)$。
\begin{questions}
\item 证明 $f$ 可积当且仅当 $\sum_i\abs{f(i)}<\infty$，此时 $\int f\dd\#=\sum_if(i)$。
\begin{solution}
由“计数测度与级数”的例子，$\int\abs f\dd\#=\sum_i\abs{f(i)}$，所以可积当且仅当级数绝对收敛。此时 $\int f\dd\#=\int f^+\dd\#-\int f^-\dd\#=\sum_if^+(i)-\sum_if^-(i)=\sum_if(i)$，最后一步是两个收敛级数可以逐项相减。
\end{solution}
\item 用逐项积分证明：对 $a_{m,n}\geqslant0$，$\sum_m\sum_na_{m,n}=\sum_n\sum_ma_{m,n}$。
\begin{solution}
令 $f_m(n)=a_{m,n}$，它是 $\N$ 上的非负函数。由逐项积分（推论~\ref{cor:integration:series}）与上一个例子，
\[
\sum_n\sum_ma_{m,n}=\int\sum_mf_m\dd\#=\sum_m\int f_m\dd\#=\sum_m\sum_na_{m,n}.
\]
数学分析中这个结论要单独证明，这里是逐项积分的特例。
\end{solution}
\end{questions}

\item\label{ex:integration:A4} 对本讲开头的例子 $f_n=\one_{\{q_1,\dots,q_n\}}$，验证 $f_n$ Riemann 可积，积分为 $0$，而极限 $\one_{\Q\cap[0,1]}$ 不是 Riemann 可积的；计算这些函数的 Lebesgue 积分，并与单调收敛定理对照。
\begin{solution}
\emph{$f_n$ 的 Riemann 积分。}任何分割的下和都是 $0$（每个小区间里都有不属于 $\{q_1,\dots,q_n\}$ 的点）。给定 $\varepsilon>0$，取分割使每个 $q_i$ 所在的小区间（至多两个）长度都小于 $\frac\varepsilon{2n}$，则上和至多 $2n\cdot\frac\varepsilon{2n}=\varepsilon$。所以上下积分都是 $0$。

\emph{极限不是 Riemann 可积的。}任何小区间都含有有理数与无理数，上和恒为 $1$，下和恒为 $0$。

\emph{Lebesgue 积分。}$\int f_n\dd\lambda=\lambda(\{q_1,\dots,q_n\})=0$，$\int\one_{\Q\cap[0,1]}\dd\lambda=\lambda(\Q\cap[0,1])=0$。$0\leqslant f_n\uparrow\one_{\Q\cap[0,1]}$，单调收敛定理说积分 $0\uparrow0$，完全一致。Riemann 积分在这里失败，Lebesgue 积分没有任何问题。
\end{solution}

\item\label{ex:integration:A5} 单调收敛定理的推广。
\begin{questions}
\item 举例说明单调收敛定理中 $f_n\geqslant0$ 的条件不能去掉。
\begin{solution}
在 $(\R,\cR,\lambda)$ 上取 $f_n=-\one_{[n,\infty)}$。$f_n\uparrow0$，但按本讲的约定 $\int f_n\dd\lambda=-\lambda([n,\infty))=-\infty$ 对一切 $n$ 成立，而 $\int0\dd\lambda=0$。问题在于“无穷多的负质量”逃到了无穷远处。
\end{solution}
\item 证明：若 $f_n\uparrow f$ 并且 $f_1$ 可积，则 $\int f_n\dd\mu\uparrow\int f\dd\mu$；若 $f_n\downarrow f$ 并且 $f_1$ 可积，则 $\int f_n\dd\mu\downarrow\int f\dd\mu$。（这里 $\int f\dd\mu$ 按本讲的约定可以是 $\pm\infty$。）
\begin{solution}
\emph{一个辅助事实。}若 $g\geqslant0$ 可测，$h$ 可积，则 $\int(g+h)\dd\mu=\int g\dd\mu+\int h\dd\mu$（两边都在 $(-\infty,\infty]$ 中）。证明：$(g+h)^-\leqslant h^-$，所以 $\int(g+h)^-<\infty$，$\int(g+h)$ 有定义。由 $(g+h)^+-(g+h)^-=g+h^+-h^-$ 得 $(g+h)^++h^-=g+h^++(g+h)^-$，四项都非负，由定理~\ref{thm:integration:linear-nonnegative} 积分后移项（有限的量可以移）即得。

\emph{递增。}$g_n=f_n-f_1\geqslant0$，$g_n\uparrow f-f_1$（$f_1$ 实值，差有意义）。由单调收敛定理 $\int g_n\uparrow\int(f-f_1)$。由辅助事实，$\int f_n=\int g_n+\int f_1$，$\int f=\int(f-f_1)+\int f_1$，所以 $\int f_n\uparrow\int f$。

\emph{递减。}$f_1-f_n\geqslant0$，$f_1-f_n\uparrow f_1-f$，由单调收敛定理 $\int(f_1-f_n)\uparrow\int(f_1-f)$。$f_n\leqslant f_1$，$\int f_n^+\leqslant\int f_1^+<\infty$，$\int f_n$ 有定义，取值于 $[-\infty,\infty)$；与上面同理 $\int f_1-\int f_n=\int(f_1-f_n)$，所以 $\int f_n\downarrow\int f$。
\end{solution}
\end{questions}

\item\label{ex:integration:A6} 设 $f$ 可积。
\begin{questions}
\item 证明 $n\to\infty$ 时 $\int_{\{\abs f>n\}}\abs f\dd\mu\to0$。
\begin{solution}
$\abs f\one_{\{\abs f\leqslant n\}}\uparrow\abs f\one_{\{\abs f<\infty\}}$，而 $\abs f<\infty$ a.e.（命题~\ref{prop:integration:markov}(3)）。由单调收敛定理，$\int_{\{\abs f\leqslant n\}}\abs f\dd\mu\uparrow\int\abs f\dd\mu<\infty$。两者之差 $\int_{\{\abs f>n\}}\abs f\dd\mu\to0$（差可以计算，因为都是有限数，由线性性）。
\end{solution}
\item 又设 $E_n\uparrow\Omega$，证明 $\int_{E_n^c}\abs f\dd\mu\to0$。
\begin{solution}
$\abs f\one_{E_n}\uparrow\abs f$，由单调收敛定理 $\int_{E_n}\abs f\uparrow\int\abs f<\infty$，所以 $\int_{E_n^c}\abs f=\int\abs f-\int_{E_n}\abs f\to0$。这两个结论说的是：可积函数的积分主要集中在“$f$ 不太大”的地方与一个“测度有限”的地方，以后讨论一致可积性时会用到。
\end{solution}
\end{questions}

\item\label{ex:integration:A7} 设 $f\geqslant0$ 可测，$A\in\cF$。
\begin{questions}
\item 证明：若 $\mu(A)=0$，则 $\int_Af\dd\mu=0$。
\begin{solution}
$f\one_A$ 在 $A^c$ 上为零，所以 $f\one_A=0$ a.e.，由命题~\ref{prop:integration:nonnegative-basic}(3)，$\int f\one_A\dd\mu=\int0\dd\mu=0$。
\end{solution}
\item 若 $f>0$ 处处成立且 $\int_Af\dd\mu=0$，证明 $\mu(A)=0$。
\begin{solution}
由命题~\ref{prop:integration:markov}(2)，$f\one_A=0$ a.e.。而 $f>0$，所以 $\{f\one_A\neq0\}=A$，$\mu(A)=0$。
\end{solution}
\end{questions}

\item\label{ex:integration:A8} 设 $V\subseteq[0,1]$ 是任意集合（不一定可测），在定理~\ref{thm:integration:bounded} 中把 $f$ 换成 $\one_V$，$E=[0,1]$，$\varphi,\psi$ 仍取 Borel 可测的简单函数。
\begin{questions}
\item 证明
\[
\begin{aligned}
\sup_{\varphi\leqslant\one_V}\int\varphi\dd\lambda&=\sup\{\lambda(A):A\subseteq V,\ A\in\cR\},\\
\inf_{\psi\geqslant\one_V}\int\psi\dd\lambda&=\inf\{\lambda(B):B\supseteq V,\ B\in\cR\}.
\end{aligned}
\]
\begin{solution}
这里 $\varphi,\psi$ 都在 $[0,1]^c$ 上为零。

\emph{第一式。}“$\geqslant$”：对 Borel 集 $A\subseteq V$，$\varphi=\one_A\leqslant\one_V$，积分是 $\lambda(A)$。“$\leqslant$”：设 $\varphi\leqslant\one_V$，令 $A=\{\varphi>0\}$，它是 Borel 集；在 $V^c$ 上 $\varphi\leqslant0$，所以 $A\subseteq V$；又 $\varphi\leqslant1$，所以 $\varphi\leqslant\one_A$，$\int\varphi\leqslant\lambda(A)$。

\emph{第二式。}“$\leqslant$”：对 Borel 集 $B\supseteq V$，$\psi=\one_{B\cap[0,1]}\geqslant\one_V$，积分是 $\lambda(B\cap[0,1])\leqslant\lambda(B)$。“$\geqslant$”：设 $\psi\geqslant\one_V$（$\psi$ 在 $[0,1]^c$ 上为零），令 $B=\{\psi\geqslant1\}$，它是 Borel 集且 $B\supseteq V$；$\psi\geqslant\one_V\geqslant0$ 处处成立，所以 $\psi\geqslant\one_B$，$\int\psi\geqslant\lambda(B)$。
\end{solution}
\item 说明：若二者相等，则存在 Borel 集 $A\subseteq V\subseteq B$ 使 $\lambda(B\setminus A)=0$。这说明定理~\ref{thm:integration:bounded} 中可测性的假设是必要的。
\begin{solution}
设公共值为 $c$。取 Borel 集 $A_k\subseteq V$，$\lambda(A_k)>c-\frac1k$，Borel 集 $B_k\supseteq V$，$\lambda(B_k)<c+\frac1k$。令 $A=\bigcup_kA_k\subseteq V$，$B=\bigcap_kB_k\cap[0,1]\supseteq V$，则 $c\geqslant\lambda(A)\geqslant c$，$c\leqslant\lambda(B)\leqslant c$，于是 $\lambda(B\setminus A)=0$。

所以 $V$ 被两个只差一个零测集的 Borel 集夹住，由第~\ref{lec:measures} 讲作业 B8)，$V$ 是 Lebesgue 可测的。反过来，若 $V$ 不是 Lebesgue 可测的（第~\ref{lec:extension-of-measures} 讲作业 D），上确界严格小于下确界，定理~\ref{thm:integration:bounded} 的等式对 $\one_V$ 不成立。这就是我们要求 $f$ 可测的原因。
\end{solution}
\end{questions}

\item\label{ex:integration:A9} 证明 $\int_0^\infty\frac{\abs{\sin x}}x\dd x=\infty$。（提示：在每个 $[k\pi,(k+1)\pi]$ 上估计，再用逐项积分。）
\begin{solution}
被积函数在 $0$ 处补充定义为 $1$，它在 $[0,\infty)$ 上连续、非负。在 $I_k=[k\pi,(k+1)\pi)$ 上，$\frac{\abs{\sin x}}x\geqslant\frac{\abs{\sin x}}{(k+1)\pi}$，而 $\int_{k\pi}^{(k+1)\pi}\abs{\sin x}\dd x=2$（Riemann 积分，由定理~\ref{thm:integration:riemann} 等于 Lebesgue 积分）。所以 $\int_{I_k}\frac{\abs{\sin x}}x\dd x\geqslant\frac2{(k+1)\pi}$。$[0,\infty)$ 是 $I_k$ 的不交并，由逐项积分（推论~\ref{cor:integration:series}），
\[
\int_0^\infty\frac{\abs{\sin x}}x\dd x=\sum_{k=0}^\infty\int_{I_k}\frac{\abs{\sin x}}x\dd x\geqslant\frac2\pi\sum_{k=0}^\infty\frac1{k+1}=\infty .
\]
\end{solution}

\item\label{ex:integration:A10} 设 $p>0$。
\begin{questions}
\item 证明 $\int_{(0,1]}x^{-p}\dd x<\infty$ 当且仅当 $p<1$。
\begin{solution}
$f_n=x^{-p}\one_{[1/n,1]}\uparrow x^{-p}\one_{(0,1]}$。$f_n$ 在 $[\frac1n,1]$ 上连续，由定理~\ref{thm:integration:riemann}，
\[
\int f_n\dd\lambda=\int_{1/n}^1x^{-p}\dd x=\begin{cases}\dfrac{1-n^{p-1}}{1-p},&p\neq1,\\[1ex] \log n,&p=1.\end{cases}
\]
由单调收敛定理，所求积分是 $n\to\infty$ 时的极限：$p<1$ 时为 $\frac1{1-p}$；$p\geqslant1$ 时为 $\infty$。
\end{solution}
\item 证明 $\int_{[1,\infty)}x^{-p}\dd x<\infty$ 当且仅当 $p>1$。
\begin{solution}
$f_n=x^{-p}\one_{[1,n]}\uparrow x^{-p}\one_{[1,\infty)}$，$\int f_n\dd\lambda=\frac{n^{1-p}-1}{1-p}$（$p\neq1$）或 $\log n$（$p=1$）。极限在 $p>1$ 时为 $\frac1{p-1}$，在 $p\leqslant1$ 时为 $\infty$。

两个结论合起来：$x^{-p}$ 在 $0$ 附近要求 $p$ 小，在无穷远处要求 $p$ 大，没有一个 $p$ 使 $x^{-p}$ 在 $(0,\infty)$ 上可积。这两个积分是以后判断矩是否存在的标准比较对象。
\end{solution}
\end{questions}

\item\label{ex:integration:A11} （反向 Fatou 引理）设 $f_n$ 可测，存在可积函数 $g$ 使 $f_n\leqslant g$ 对一切 $n$ 成立。证明
\[
\limsup_n\int f_n\dd\mu\leqslant\int\limsup_nf_n\dd\mu .
\]
（这里 $\int\limsup_nf_n\dd\mu$ 可以是 $-\infty$。）
\begin{solution}
思路：把“上方有界”化为“非负”，用 Fatou 引理。

令 $h_n=g-f_n$。由 $f_n\leqslant g$，$h_n\geqslant0$；它可测（$g$ 可积从而 a.e. 有限，像控制收敛定理证明第一步那样，在零测集上把 $g$ 改成 $0$、$f_n$ 改成 $0$ 不影响任何积分，所以不妨设 $g$ 处处有限）。由 Fatou 引理（定理~\ref{thm:integration:fatou}），
\[
\int\liminf_n(g-f_n)\dd\mu\leqslant\liminf_n\int(g-f_n)\dd\mu .
\]
左边：$\liminf_n(g-f_n)=g-\limsup_nf_n$。右边：由于 $f_n^+\leqslant g^+$，$\int f_n^+\dd\mu<\infty$，所以 $\int f_n\dd\mu$ 按本讲的约定有意义，取值于 $[-\infty,\infty)$。我们说明 $\int(g-f_n)\dd\mu=\int g\dd\mu-\int f_n\dd\mu$：当 $\int f_n$ 有限时，这是线性性（定理~\ref{thm:integration:integrable}(3)）；当 $\int f_n=-\infty$，即 $\int f_n^-=\infty$ 时，由 $g-f_n=g-f_n^++f_n^-\geqslant f_n^--2\abs g$ 得 $\int(g-f_n)\geqslant\int f_n^--2\int\abs g=\infty$，两边都是 $+\infty$。于是右边等于 $\int g\dd\mu-\limsup_n\int f_n\dd\mu$。

同样地，$\limsup_nf_n\leqslant g$，所以 $\int\limsup_nf_n\dd\mu$ 有意义（可以是 $-\infty$），左边等于 $\int g\dd\mu-\int\limsup_nf_n\dd\mu$。得
\[
\int g\dd\mu-\int\limsup_nf_n\dd\mu\leqslant\int g\dd\mu-\limsup_n\int f_n\dd\mu ,
\]
消去有限数 $\int g\dd\mu$ 即得结论。

注意：控制收敛定理正是 Fatou 引理与反向 Fatou 引理的结合：$\int\liminf f_n\leqslant\liminf\int f_n\leqslant\limsup\int f_n\leqslant\int\limsup f_n$，当 $f_n\to f$ 时两端都等于 $\int f$。
\end{solution}

\item\label{ex:integration:A12} 计算 $\displaystyle\lim_{n\to\infty}\int_0^\infty\frac{n\sin(x/n)}{x(1+x^2)}\dd x$。
\begin{solution}
思路：被积函数逐点收敛，找一个控制函数。

令 $f_n(x)=\frac{n\sin(x/n)}{x(1+x^2)}$，$x>0$。它在 $(0,\infty)$ 上连续，是 Borel 函数。

\emph{逐点极限。}固定 $x>0$，由 $\lim_{y\to0}\frac{\sin y}y=1$，$n\sin(x/n)=x\cdot\frac{\sin(x/n)}{x/n}\to x$，所以 $f_n(x)\to\frac1{1+x^2}$。

\emph{控制函数。}对一切实数 $y$，$\abs{\sin y}\leqslant\abs y$。所以 $\abs{n\sin(x/n)}\leqslant n\cdot\frac xn=x$，从而 $\abs{f_n(x)}\leqslant\frac1{1+x^2}=:g(x)$。$g$ 可积：由单调收敛定理，$\int_{(0,\infty)}g\dd\lambda=\lim_N\arctan N=\frac\pi2$。

由控制收敛定理，极限等于 $\int_0^\infty\frac{\dd x}{1+x^2}=\frac\pi2$。
\end{solution}

\item\label{ex:integration:A13} （广义控制收敛定理）设 $f_n\to f$ a.e.，$g_n\to g$ a.e.，$\abs{f_n}\leqslant g_n$，$g_n,g$ 可积，并且 $\int g_n\dd\mu\to\int g\dd\mu$。证明 $\int f_n\dd\mu\to\int f\dd\mu$。
\begin{solution}
思路：和控制收敛定理的证明一样，对两个非负函数列 $g_n+f_n$ 与 $g_n-f_n$ 用 Fatou 引理。

和控制收敛定理证明的第一步一样，去掉一个零测集之后不妨设处处收敛，$g_n,g$ 处处有限。令 $n\to\infty$，$\abs f\leqslant g$，所以 $f$ 可积。

$g_n+f_n\geqslant0$，并且 $g_n+f_n\to g+f$。由 Fatou 引理与线性性，
\[
\int g\dd\mu+\int f\dd\mu\leqslant\liminf_n\Bigl(\int g_n\dd\mu+\int f_n\dd\mu\Bigr)=\int g\dd\mu+\liminf_n\int f_n\dd\mu ,
\]
这里用到了 $\int g_n\to\int g$（一个收敛数列与另一个数列之和的下极限，等于极限加下极限）。消去 $\int g$（有限），得 $\int f\leqslant\liminf_n\int f_n$。

$g_n-f_n\geqslant0$，$g_n-f_n\to g-f$。同样得到 $\int g-\int f\leqslant\int g-\limsup_n\int f_n$，即 $\limsup_n\int f_n\leqslant\int f$。

两式合起来，$\lim_n\int f_n\dd\mu=\int f\dd\mu$。控制收敛定理是 $g_n=g$ 的特例。
\end{solution}

\item\label{ex:integration:A14} （Scheffé 引理）设 $f_n,f\geqslant0$ 可积，$f_n\to f$ a.e.，并且 $\int f_n\dd\mu\to\int f\dd\mu$。
\begin{questions}
\item 证明 $\int(f-f_n)^+\dd\mu\to0$。
\begin{solution}
思路：$f-f_n$ 的正部被 $f$ 控制，可以用控制收敛定理；负部用总积分的条件来处理。

记 $(f-f_n)^+=\max\{f-f_n,0\}$，$(f-f_n)^-=\max\{f_n-f,0\}$。

由于 $f_n\geqslant0$，$f-f_n\leqslant f$，所以 $0\leqslant(f-f_n)^+\leqslant f$。又 $(f-f_n)^+\to0$ a.e.（$x\mapsto\max\{x,0\}$ 连续）。$f$ 可积，由控制收敛定理，$\int(f-f_n)^+\dd\mu\to0$。
\end{solution}
\item 证明 $\int\abs{f_n-f}\dd\mu\to0$。特别地，若 $f_n,f$ 都是概率密度，并且 $f_n\to f$ a.e.，则 $\int\abs{f_n-f}\dd x\to0$。
\begin{solution}
记号同 A14-1)。由 A14-1)，正部的积分趋于零；下面处理负部。$f-f_n=(f-f_n)^+-(f-f_n)^-$，积分得
\[
\begin{aligned}
\int(f-f_n)^-\dd\mu&=\int(f-f_n)^+\dd\mu-\int f\dd\mu+\int f_n\dd\mu\\
&\to0-\int f\dd\mu+\int f\dd\mu=0 .
\end{aligned}
\]

于是 $\int\abs{f_n-f}\dd\mu=\int(f-f_n)^+\dd\mu+\int(f-f_n)^-\dd\mu\to0$。

这个结论在概率论中的意义是：若随机变量 $X_n$ 的密度 $f_n$ 几乎处处收敛到 $X$ 的密度 $f$，则对一切 Borel 集 $B$，$\abs{P(X_n\in B)-P(X\in B)}=\abs{\int_B(f_n-f)\dd x}\leqslant\int\abs{f_n-f}\dd x\to0$，并且这个收敛关于 $B$ 是一致的。
\end{solution}
\end{questions}

\item\label{ex:integration:A15} （积分的绝对连续性）设 $g$ 可积。证明：对任意 $\varepsilon>0$，存在 $\delta>0$，使得只要 $A\in\cF$，$\mu(A)<\delta$，就有 $\int_A\abs g\dd\mu<\varepsilon$。
\begin{solution}
思路：把 $\abs g$ 拆成有界的部分和“高处”的部分。有界部分在小集合上的积分小，因为它不超过界乘以集合的测度；高处的部分整体积分小，由单调收敛定理。

令 $g_M=\abs g\one_{\{\abs g\leqslant M\}}$。当 $M\uparrow\infty$（取 $M=1,2,\dots$）时，$g_M\uparrow\abs g\one_{\{\abs g<\infty\}}$，而 $\abs g<\infty$ a.e.，所以由单调收敛定理，$\int g_M\dd\mu\uparrow\int\abs g\dd\mu$。$\int\abs g$ 有限，所以
\[
\int\abs g\one_{\{\abs g>M\}}\dd\mu=\int\abs g\dd\mu-\int g_M\dd\mu\to0 .
\]
给定 $\varepsilon>0$，取整数 $M$ 使 $\int\abs g\one_{\{\abs g>M\}}\dd\mu<\varepsilon/2$，再取 $\delta=\frac\varepsilon{2M}$。若 $\mu(A)<\delta$，则
\[
\begin{aligned}
\int_A\abs g\dd\mu&=\int_A\abs g\one_{\{\abs g\leqslant M\}}\dd\mu+\int_A\abs g\one_{\{\abs g>M\}}\dd\mu\\
&\leqslant M\mu(A)+\int\abs g\one_{\{\abs g>M\}}\dd\mu<\frac\varepsilon2+\frac\varepsilon2=\varepsilon .
\end{aligned}
\]
（第一项的估计：$\abs g\one_{\{\abs g\leqslant M\}}\one_A\leqslant M\one_A$，由单调性与 $\int M\one_A=M\mu(A)$。）

下一讲讨论测度的绝对连续性时，我们会从另一个角度再看到这个结论：$\nu(A)=\int_A\abs g\dd\mu$ 是有限测度，而上面的结论说 $\mu(A)$ 小时 $\nu(A)$ 也小。
\end{solution}

\item\label{ex:integration:A16} 设 $f$ 在 $[a,b]$ 上关于 Lebesgue 测度可积，$g(x)=\int_{[a,x]}f\dd\lambda$（$a\leqslant x\leqslant b$）。证明 $g$ 在 $[a,b]$ 上连续。
\begin{solution}
思路：$g(y)-g(x)$ 是 $f$ 在一个短区间上的积分，区间的长度趋于零时，由习题 A15，积分趋于零。

设 $a\leqslant x<y\leqslant b$。由于 $[a,y]=[a,x]+(x,y]$，由积分对积分区域的可加性（$\int_{A+B}f=\int_Af+\int_Bf$：对 $f^+$ 与 $f^-$ 分别用推论~\ref{cor:integration:indefinite}，再相减），
\[
\abs{g(y)-g(x)}=\Bigl\lvert\int_{(x,y]}f\dd\lambda\Bigr\rvert\leqslant\int_{(x,y]}\abs f\dd\lambda .
\]
给定 $\varepsilon>0$，对可积函数 $\abs f\one_{[a,b]}$ 用习题 A15，取 $\delta>0$，使 $\lambda(A)<\delta$ 时 $\int_A\abs f\one_{[a,b]}\dd\lambda<\varepsilon$。那么只要 $\abs{y-x}<\delta$，$\lambda((x,y])=y-x<\delta$，就有 $\abs{g(y)-g(x)}<\varepsilon$。所以 $g$ 在 $[a,b]$ 上一致连续，当然连续。

也可以用控制收敛定理证明：设 $x_n\to x$，则对 $t\neq x$，$\one_{[a,x_n]}(t)\to\one_{[a,x]}(t)$，即 $f\one_{[a,x_n]}\to f\one_{[a,x]}$ 几乎处处成立（例外的点只有 $t=x$，它是零测集），并且 $\abs{f\one_{[a,x_n]}}\leqslant\abs f\one_{[a,b]}$，由控制收敛定理（定理~\ref{thm:integration:dct}），$g(x_n)\to g(x)$。注意在 $t=x$ 处收敛可能不成立（$x_n\uparrow x$，$x_n<x$ 时 $\one_{[a,x_n]}(x)=0$），所以这里“几乎处处”是必要的。
\end{solution}
\end{exercises}


\begin{probleme}{B}{Riemann 可积性的 Lebesgue 判别法}
设 $f:[a,b]\to\R$ 有界，$\abs f\leqslant M$，不假设可测。对 $x\in[a,b]$ 与 $\delta>0$ 令
\[
f^\delta(x)=\sup\{f(y):y\in[a,b],\ \abs{y-x}<\delta\},\qquad f_\delta(x)=\inf\{f(y):y\in[a,b],\ \abs{y-x}<\delta\},
\]
$f^0(x)=\lim_{\delta\downarrow0}f^\delta(x)$，$f_0(x)=\lim_{\delta\downarrow0}f_\delta(x)$（单调极限）。本题证明：\emph{$f$ Riemann 可积当且仅当 $f$ 的间断点集是 Lebesgue 零测集。}
\begin{questions}
\item 证明 $f^0$ 上半连续：对任意 $c$，$\{x:f^0(x)<c\}$ 是 $[a,b]$ 中的相对开集。同理 $f_0$ 下半连续。从而 $f^0,f_0$ 是 Borel 函数，并且 $f_0\leqslant f\leqslant f^0$。
\begin{solution}
设 $f^0(x)<c$。由于 $f^\delta(x)\downarrow f^0(x)$，存在 $\delta>0$ 使 $f^\delta(x)<c$。若 $x'\in[a,b]$，$\abs{x'-x}<\delta/2$，则以 $x'$ 为心、$\delta/2$ 为半径的区间含于以 $x$ 为心、$\delta$ 为半径的区间，所以 $f^0(x')\leqslant f^{\delta/2}(x')\leqslant f^\delta(x)<c$。于是 $\{f^0<c\}$ 包含 $x$ 的一个相对邻域，是相对开集，也就是某个开集与 $[a,b]$ 的交，是 Borel 集。由推论~\ref{cor:random-variables:criterion}（用 $\{f^0<c\}$ 的形式），$f^0$ 是 Borel 函数。$f_0$ 同理（考虑 $\{f_0>c\}$）。$y=x$ 在 $x$ 的每个邻域里，所以 $f_\delta(x)\leqslant f(x)\leqslant f^\delta(x)$，取极限得 $f_0\leqslant f\leqslant f^0$。
\end{solution}
\item 证明 $f$ 在 $x$ 处连续当且仅当 $f^0(x)=f_0(x)$。
\begin{solution}
与第~\ref{lec:random-variables} 讲作业 C4) 相同。若 $f$ 在 $x$ 连续，对 $\varepsilon>0$ 取 $\delta$ 使邻域内 $\abs{f(y)-f(x)}<\varepsilon$，则 $f^\delta(x)-f_\delta(x)\leqslant2\varepsilon$，从而 $f^0(x)=f_0(x)$。反之，若 $f^0(x)=f_0(x)$，二者都等于 $f(x)$；对 $\varepsilon>0$ 取 $\delta$ 使 $f^\delta(x)<f(x)+\varepsilon$，$f_\delta(x)>f(x)-\varepsilon$，则邻域内 $\abs{f(y)-f(x)}<\varepsilon$。
\end{solution}
\item 令 $P_k$ 为把 $[a,b]$ 等分成 $2^k$ 份的分割，$s_{P_k},t_{P_k}$ 如定理~\ref{thm:integration:riemann} 的证明。证明 $s_{P_k}$ 单调不减，$t_{P_k}$ 单调不增，并且在不属于任何 $P_k$ 分点的 $x$ 处，$t_{P_k}(x)\to f^0(x)$，$s_{P_k}(x)\to f_0(x)$。
\begin{solution}
\emph{单调。}$P_{k+1}$ 把 $P_k$ 的每个小区间一分为二。设 $x$ 在 $P_{k+1}$ 的小区间 $J$ 中，$J$ 含于 $P_k$ 的小区间 $I$。$t_{P_{k+1}}(x)=\sup_Jf\leqslant\sup_If=t_{P_k}(x)$；$s$ 同理。（这里 $J$、$I$ 指闭区间，与定理证明中 $M_k,m_k$ 的定义一致。）

\emph{极限。}设 $x$ 不是任何 $P_k$ 的分点。记 $I_k(x)$ 为 $P_k$ 中含 $x$ 的闭小区间，长度 $(b-a)2^{-k}$，$x$ 在它的内部。一方面，给定 $\delta>0$，当 $(b-a)2^{-k}<\delta$ 时 $I_k(x)\subseteq(x-\delta,x+\delta)$，所以 $t_{P_k}(x)=\sup_{I_k(x)}f\leqslant f^\delta(x)$。另一方面，$x$ 是 $I_k(x)$ 的内点，存在 $\delta_k>0$ 使 $(x-\delta_k,x+\delta_k)\subseteq I_k(x)$，所以 $t_{P_k}(x)\geqslant f^{\delta_k}(x)\geqslant f^0(x)$。于是 $f^0(x)\leqslant\lim_kt_{P_k}(x)\leqslant f^\delta(x)$ 对一切 $\delta$ 成立，令 $\delta\downarrow0$ 得 $\lim_kt_{P_k}(x)=f^0(x)$。$s$ 同理。
\end{solution}
\item 由于全体分点是可数集，B3) 中的收敛几乎处处成立。对 $M-t_{P_k}\geqslant0$ 与 $s_{P_k}+M\geqslant0$ 用单调收敛定理，证明 $U(P_k)\to\int_a^bf^0\dd\lambda$，$L(P_k)\to\int_a^bf_0\dd\lambda$。
\begin{solution}
所有分点组成可数集，Lebesgue 测度为零。在 $[a,b]$ 上，$M-t_{P_k}$ 非负、单调不减，几乎处处趋于 $M-f^0$。把它们乘以 $\one_{[a,b]}$，由单调收敛定理（在零测集上修改函数不影响积分），
\[
M(b-a)-U(P_k)=\int_{[a,b]}(M-t_{P_k})\dd\lambda\to\int_{[a,b]}(M-f^0)\dd\lambda=M(b-a)-\int_{[a,b]}f^0\dd\lambda .
\]
这里用了 $\int_{[a,b]}t_{P_k}\dd\lambda=U(P_k)$（定理~\ref{thm:integration:riemann} 的证明）以及各项有限，可以移项。所以 $U(P_k)\to\int_a^bf^0\dd\lambda$。下和同理。
\end{solution}
\item 利用数学分析中的 Darboux 定理（分割的细度趋于零时，上、下和分别收敛到上、下积分），证明 $f$ Riemann 可积当且仅当 $\int_a^b(f^0-f_0)\dd\lambda=0$，再由命题~\ref{prop:integration:markov}(2) 与 B2) 完成证明。
\begin{solution}
由 Darboux 定理，上积分 $=\lim_kU(P_k)=\int_a^bf^0\dd\lambda$，下积分 $=\lim_kL(P_k)=\int_a^bf_0\dd\lambda$。$f$ Riemann 可积当且仅当二者相等，即 $\int_a^b(f^0-f_0)\dd\lambda=0$。$f^0-f_0\geqslant0$，由命题~\ref{prop:integration:markov}(2)，这等价于 $f^0=f_0$ a.e. 于 $[a,b]$，由 B2)，又等价于间断点集 $\{f^0\neq f_0\}$ 是零测集。
\end{solution}
\item 证明 Riemann 可积的 $f$ 是 Lebesgue 可测的（对 Lebesgue 可测集的 $\sigma$-代数），并且它的 Lebesgue 积分等于 Riemann 积分。（提示：$f_0\leqslant f\leqslant f^0$，且 $f^0=f_0$ a.e.；用第~\ref{lec:extension-of-measures} 讲的完备性。）
\begin{solution}
设 $f$ Riemann 可积，$N=\{f_0<f^0\}$ 是 Borel 零测集，在 $N^c$ 上 $f=f^0$。对任意 $c$，
\[
\{f\leqslant c\}=\bigl(\{f^0\leqslant c\}\setminus N\bigr)\cup\bigl(\{f\leqslant c\}\cap N\bigr).
\]
第一部分是 Borel 集，第二部分是零测集 $N$ 的子集，由 Carathéodory 定理（定理~\ref{thm:extension:caratheodory} 的最后一句）是 Lebesgue 可测的。所以 $f$ 是 Lebesgue 可测的。在完备化后的测度下，$f=f^0$ a.e.，所以 $\int_{[a,b]}f\dd\lambda=\int_a^bf^0\dd\lambda$，由 B5) 的论证，它等于上积分，即 Riemann 积分。

这去掉了定理~\ref{thm:integration:riemann} 中“Borel 可测”的假设。
\end{solution}
\item 用这个判别法说明：Riemann 函数（在既约分数 $p/q$ 处取 $1/q$，在无理数处取 $0$）Riemann 可积；Cantor 集的示性函数 $\one_C$ Riemann 可积；$\one_{\Q\cap[0,1]}$ 不是。
\begin{solution}
\emph{Riemann 函数 $R$。}在无理数 $x$ 处连续：给定 $\varepsilon>0$，$[0,1]$ 中分母 $q\leqslant1/\varepsilon$ 的有理数只有有限个，取 $\delta$ 使 $(x-\delta,x+\delta)$ 不含它们，则邻域内 $R<\varepsilon$，而 $R(x)=0$。在有理数 $p/q$ 处不连续：$R(p/q)=1/q$，而附近的无理数处 $R=0$。间断点集是 $\Q\cap[0,1]$，零测集，所以 Riemann 可积（积分为 $0$）。

\emph{$\one_C$。}$[0,1]\setminus C$ 是开集，在其中 $\one_C=0$ 局部为常数，连续。在 $x\in C$ 处，$\one_C(x)=1$；$C$ 不含任何区间（它的测度为 $0$），所以 $x$ 的任何邻域中都有 $C^c$ 的点，$\one_C$ 在 $x$ 处不连续。间断点集是 $C$，零测集，所以 Riemann 可积（积分为 $0$）。

\emph{$\one_{\Q\cap[0,1]}$。}处处不连续，间断点集 $[0,1]$ 的测度为 $1$，不是 Riemann 可积的。
\end{solution}
\end{questions}
\end{probleme}

\begin{probleme}{C}{分层积分公式}
设 $f:\Omega\to[0,\infty]$ 可测。本题证明
\[
\int f\dd\mu=\int_{(0,\infty)}\mu(\{f>t\})\dd t,
\]
其中右边是 $(0,\infty)$ 上关于 Lebesgue 测度的积分。直观上，左边按“竖条”计算 $f$ 的图像下方的面积，右边按“横条”计算：高度 $t$ 处横条的长度是 $\mu(\{f>t\})$。这正是 Lebesgue 分割值域的想法。
\begin{questions}
\item 令 $G(t)=\mu(\{f>t\})$，$t>0$。证明 $G$ 单调不增、右连续，从而是 Borel 函数（取值于 $[0,\infty]$ 的单调函数同样可测）。
\begin{solution}
$s<t$ 时 $\{f>t\}\subseteq\{f>s\}$，所以 $G$ 单调不增。若 $t_n\downarrow t$，则 $\{f>t_n\}\uparrow\{f>t\}$（$f(\omega)>t$ 时，$n$ 充分大有 $f(\omega)>t_n$），由从下连续 $G(t_n)\uparrow G(t)$，所以 $G$ 右连续（这一步不需要有限性）。可测性：对实数 $c$，$\{t>0:G(t)\leqslant c\}$ 具有性质“若 $t$ 属于它且 $s>t$，则 $s$ 也属于它”，所以是一个区间，是 Borel 集；由命题~\ref{prop:random-variables:extended}，$G$ 是广义实值 Borel 函数。
\end{solution}
\item 设 $f=\sum_{j=1}^my_j\one_{A_j}$，$0<y_1<\cdots<y_m$，$A_j$ 两两不交。写出 $G$，并验证公式成立（注意某些 $\mu(A_j)$ 可能是 $\infty$）。
\begin{solution}
$f(\omega)>t$ 当且仅当 $\omega\in A_j$ 且 $y_j>t$。所以 $G(t)=\sum_{j:\,y_j>t}\mu(A_j)=\sum_j\mu(A_j)\one_{(0,y_j)}(t)$。这是 $m$ 个非负函数之和，第 $j$ 个是常数 $\mu(A_j)$（可能为 $\infty$）乘以 $\one_{(0,y_j)}$，它的积分是 $\mu(A_j)\cdot y_j$（$\mu(A_j)=\infty$ 时，$n\one_{(0,y_j)}\uparrow\infty\cdot\one_{(0,y_j)}$，由单调收敛定理积分为 $\infty$）。由线性性，
\[
\int_{(0,\infty)}G\dd t=\sum_jy_j\mu(A_j)=\int f\dd\mu,
\]
最后一步是推论~\ref{cor:integration:nonnegative-simple}。
\end{solution}
\item 对一般的 $f$，取非负简单函数 $\varphi_n\uparrow f$。证明对每个 $t>0$，$\{\varphi_n>t\}\uparrow\{f>t\}$，从而 $\mu(\{\varphi_n>t\})\uparrow\mu(\{f>t\})$。
\begin{solution}
$\varphi_n$ 递增，所以 $\{\varphi_n>t\}$ 递增；$\varphi_n\leqslant f$，所以 $\{\varphi_n>t\}\subseteq\{f>t\}$。若 $f(\omega)>t$，由 $\varphi_n(\omega)\uparrow f(\omega)$，$n$ 充分大时 $\varphi_n(\omega)>t$。所以 $\bigcup_n\{\varphi_n>t\}=\{f>t\}$。由从下连续得第二个结论。
\end{solution}
\item 在 $\Omega$ 上和在 $(0,\infty)$ 上各用一次单调收敛定理，完成证明。
\begin{solution}
记 $G_n(t)=\mu(\{\varphi_n>t\})$。由 C2)（每个 $\varphi_n$ 写成标准表示，去掉值为 $0$ 的部分），$\int\varphi_n\dd\mu=\int_{(0,\infty)}G_n\dd t$。在 $\Omega$ 上用单调收敛定理，左边 $\uparrow\int f\dd\mu$；由 C3)，$G_n\uparrow G$ 逐点，在 $(0,\infty)$ 上用单调收敛定理，右边 $\uparrow\int G\dd t$。所以两者相等。
\end{solution}
\item 设 $X$ 是概率空间上取非负整数值的随机变量。证明 $\int X\dd P=\sum_{k=1}^\infty P(X\geqslant k)$。
\begin{solution}
对 $t\in[k-1,k)$（$k\geqslant1$），$X>t$ 当且仅当 $X\geqslant k$（$X$ 取整数值），所以 $G(t)=P(X\geqslant k)$。$(0,\infty)$ 是 $(0,1)$ 与 $[k-1,k)$（$k\geqslant2$）的不交并，由逐项积分，$\int_{(0,\infty)}G\dd t=\sum_{k\geqslant1}P(X\geqslant k)\cdot1$。再由 C4)。这就是第~\ref{lec:expectation} 讲中尾概率公式的离散形式。
\end{solution}
\item 设 $p>0$。证明 $\int f^p\dd\mu=\int_{(0,\infty)}pt^{p-1}\mu(\{f>t\})\dd t$。（提示：对 $f^p$ 用公式，再在每个有界区间 $[\varepsilon,N]$ 上用 Riemann 积分的换元公式与定理~\ref{thm:integration:riemann}，单调函数是 Riemann 可积的，最后用单调收敛定理令 $\varepsilon\to0$，$N\to\infty$。）
\begin{solution}
对 $f^p$ 用公式：$\int f^p\dd\mu=\int_{(0,\infty)}\mu(\{f^p>s\})\dd s=\int_{(0,\infty)}G(s^{1/p})\dd s$。

\emph{情形一：某个 $t_0>0$ 使 $G(t_0)=\infty$。}则左边 $\geqslant t_0^p\mu(\{f>t_0\})=\infty$；右边，$G$ 单调不增，在 $(0,t_0]$ 上为 $\infty$，积分也是 $\infty$。

\emph{情形二：$G$ 处处有限。}在 $[\varepsilon,N]$ 上，$G$ 单调有界，是 Riemann 可积的；$s=t^p$ 是 $C^1$ 的严格递增变换。由 Riemann 积分的换元公式，
\[
\int_{\varepsilon^p}^{N^p}G(s^{1/p})\dd s=\int_\varepsilon^NG(t)\,pt^{p-1}\dd t .
\]
两个被积函数都是有界 Borel 函数，由定理~\ref{thm:integration:riemann}，两边也是 Lebesgue 积分。令 $\varepsilon\downarrow0$，$N\uparrow\infty$，被积函数乘以区间的示性函数后单调递增，由单调收敛定理两边分别趋于 $\int_{(0,\infty)}G(s^{1/p})\dd s$ 与 $\int_{(0,\infty)}pt^{p-1}G(t)\dd t$。

概率中最常用的是 $p=1,2$：$EX=\int_0^\infty P(X>t)\dd t$，$EX^2=\int_0^\infty2tP(X>t)\dd t$（$X\geqslant0$）。第~\ref{lec:product-measures} 讲会用 Fubini 定理给出另一个证明。
\end{solution}
\end{questions}
\end{probleme}

\thinkingpart

\begin{probleme}[问题]{D}{为什么要求 \texorpdfstring{$\sigma$}{σ}-有限}
在 $(\R,\cR)$ 上令 $\mu(A)=\infty$（$A\neq\varnothing$），$\mu(\varnothing)=0$。
\begin{questions}
\item 证明 $\mu$ 是测度，但不是 $\sigma$-有限的。
\begin{solution}
设 $A_n$ 两两不交。若全为空集，两边都是 $0$；否则并集非空，左边为 $\infty$，右边至少有一项为 $\infty$。所以 $\mu$ 可数可加。测度有限的集合只有 $\varnothing$，可数个空集的并不能是 $\R$，所以不是 $\sigma$-有限的。
\end{solution}
\item 若照搬定义~\ref{def:integration:nonnegative}，证明 $\int\one_\R\dd\mu=0$，而 $\mu(\R)=\infty$。
\begin{solution}
定义中的 $h$ 要求 $\mu(\{h>0\})<\infty$，只能 $\{h>0\}=\varnothing$，即 $h\equiv0$。所以对任何 $f$，“积分”都是 $0$，包括 $\int\one_\R\dd\mu=0\neq\mu(\R)$。定义~\ref{def:integration:nonnegative} 只用“测度有限的部分”去逼近，在没有任何测度有限的非空集合时就失效了。
\end{solution}
\item 另一种定义是 $\int f\dd\mu=\sup\{\sum_jy_j\mu(\{\varphi=y_j\}):0\leqslant\varphi\leqslant f,\ \varphi\ \text{简单}\}$。证明对 $\sigma$-有限的 $\mu$，两种定义给出同样的值（用推论~\ref{cor:integration:nonnegative-simple} 与单调收敛定理）。
\begin{solution}
记新定义为 $J(f)$。若 $0\leqslant\varphi\leqslant f$ 是简单函数，由推论~\ref{cor:integration:nonnegative-simple}，$\sum_jy_j\mu(\{\varphi=y_j\})=\int\varphi\dd\mu\leqslant\int f\dd\mu$，所以 $J(f)\leqslant\int f\dd\mu$。反之，取简单函数 $\varphi_n\uparrow f$，由单调收敛定理 $\int f\dd\mu=\lim_n\int\varphi_n\dd\mu\leqslant J(f)$。
\end{solution}
\item 对一般的测度，用 D3) 的定义重新证明单调收敛定理与线性性。哪一步需要修改？
\begin{solution}
\emph{单调收敛定理。}$J$ 显然单调，所以 $\lim_nJ(f_n)\leqslant J(f)$。反之，设 $0\leqslant\varphi\leqslant f$ 简单，$\varphi=\sum_jy_j\one_{\{\varphi=y_j\}}$，$c\in(0,1)$，$A_n=\{f_n\geqslant c\varphi\}\uparrow\Omega$（与定理~\ref{thm:integration:mct} 的证明相同）。$c\varphi\one_{A_n}\leqslant f_n$ 是简单函数，所以
\[
J(f_n)\geqslant c\sum_jy_j\mu(\{\varphi=y_j\}\cap A_n)\to c\sum_jy_j\mu(\{\varphi=y_j\}),
\]
这里用的是\emph{从下连续}（$\{\varphi=y_j\}\cap A_n\uparrow\{\varphi=y_j\}$），它不需要任何有限性。令 $c\uparrow1$，再对 $\varphi$ 取上确界即可。

\emph{线性性。}对非负简单函数，命题~\ref{prop:integration:simple}(3) 的计算（把两个函数写在公共的划分上）只用到有限可加性，在 $[0,\infty]$ 中没有减法，所以不需要有限性。再用单调收敛定理推到一般的非负函数，与定理~\ref{thm:integration:linear-nonnegative} 的第二步相同。

需要修改的是：正文中凡是用到“测度有限的集合 $E_n\uparrow\Omega$”与“从上连续”的地方（截断引理~\ref{lem:integration:truncation}，以及定理~\ref{thm:integration:mct} 证明中的 $\mu(S\cap A_n^c)\to0$），都换成直接对简单函数用“从下连续”。正文之所以走有界函数这一步，是为了和 Riemann 积分的 Darboux 和对照；对一般测度，D3) 的定义更直接。
\end{solution}
\end{questions}
\end{probleme}

\begin{probleme}[问题]{E}{Lebesgue 积分的平移与伸缩}
设 $f$ 是 $\R$ 上关于 $\lambda$ 可积的函数，$c\in\R$，$a\neq0$。用四步法（示性函数、简单函数、非负函数、可积函数）以及 Lebesgue 测度的平移不变性与伸缩性（第~\ref{lec:extension-of-measures} 讲），证明
\[
\int f(x+c)\dd x=\int f(x)\dd x,\qquad\int f(ax)\dd x=\frac1{\abs a}\int f(x)\dd x .
\]
在四步法的每一步中，你分别用到了什么？
\begin{solution}
只写平移；伸缩完全一样，把 $\lambda(B-c)=\lambda(B)$ 换成 $\lambda(a^{-1}B)=\abs a^{-1}\lambda(B)$（习题~\exref{ex:extension:A3}）。

\emph{第一步：示性函数。}$f=\one_B$，$B\in\cR$。$\one_B(x+c)=\one_{B-c}(x)$，$B-c$ 是 Borel 集（第~\ref{lec:set-systems} 讲的例题，第~\pageref{we:set-systems:translation} 页），所以 $\int\one_B(x+c)\dd x=\lambda(B-c)=\lambda(B)$。用到的是\emph{平移不变性}（命题~\ref{prop:extension:lebesgue}(2)）。

\emph{第二步：非负简单函数。}$f=\sum_jy_j\one_{B_j}$，$f(\cdot+c)=\sum_jy_j\one_{B_j-c}$，用第一步与\emph{线性性}。

\emph{第三步：非负可测函数。}取非负简单函数 $\varphi_n\uparrow f$，则 $\varphi_n(\cdot+c)$ 是非负简单函数，$\varphi_n(\cdot+c)\uparrow f(\cdot+c)$。两边各用一次\emph{单调收敛定理}，再用第二步。

\emph{第四步：可积函数。}$(f(\cdot+c))^\pm=f^\pm(\cdot+c)$，由第三步 $\int f^\pm(x+c)\dd x=\int f^\pm(x)\dd x<\infty$，所以 $f(\cdot+c)$ 可积，并由\emph{积分的定义} $\int f=\int f^+-\int f^-$ 得结论。
\end{solution}
\end{probleme}

\exercisepart{课外知识：Riemann--Lebesgue 引理}

\begin{probleme}{F}{Riemann--Lebesgue 引理}
设 $g$ 是 $\R$ 上关于 $\lambda$ 可积的函数。本题证明
\[
\lim_{n\to\infty}\int g(x)\cos nx\dd x=0 .
\]
这个结论在 Fourier 分析中很基本；以后讨论特征函数时，我们会看到它的一个概率论版本：有密度的分布，其特征函数在无穷远处趋于零。
\begin{questions}
\item 设 $g\geqslant0$。由定理~\ref{thm:random-variables:simple-approx} 取简单函数 $0\leqslant\varphi_k\uparrow g$，证明 $\lambda(\{\varphi_k>0\})<\infty$（提示：Markov 不等式），并由单调收敛定理证明 $\int(g-\varphi_k)\dd x\to0$。由此证明对一般可积的 $g$ 与 $\varepsilon>0$，存在 $\varphi=\sum_{i=1}^mb_i\one_{A_i}\in\cE_0$ 使 $\int\abs{g-\varphi}\dd x<\varepsilon$。
\begin{solution}
设 $c_k>0$ 是 $\varphi_k$ 取到的最小正值（若 $\varphi_k\equiv0$ 则无需证明）。$\{\varphi_k>0\}=\{\varphi_k\geqslant c_k\}$，由 Markov 不等式，$\lambda(\{\varphi_k>0\})\leqslant\frac1{c_k}\int\varphi_k\leqslant\frac1{c_k}\int g<\infty$。由单调收敛定理 $\int\varphi_k\uparrow\int g<\infty$，所以 $\int(g-\varphi_k)=\int g-\int\varphi_k\to0$。

一般的 $g$：对 $g^+$、$g^-$ 分别取 $\varphi,\psi$ 使 $\int(g^+-\varphi)<\frac\varepsilon2$，$\int(g^--\psi)<\frac\varepsilon2$。$\varphi-\psi\in\cE_0$，$\abs{g-(\varphi-\psi)}\leqslant(g^+-\varphi)+(g^--\psi)$，积分小于 $\varepsilon$。把 $\varphi-\psi$ 写成 $\sum_ib_i\one_{A_i}$（$A_i$ 两两不交，$b_i\neq0$ 时 $\lambda(A_i)<\infty$）。
\end{solution}
\item 利用第~\ref{lec:extension-of-measures} 讲作业 C7，把每个 $A_i$ 换成有限个区间的并 $B_i$，使 $\lambda(A_i\triangle B_i)$ 很小。证明存在\term{阶梯函数} $q=\sum_{j=1}^kc_j\one_{(a_{j-1},a_j]}$ 使 $\int\abs{g-q}\dd x<2\varepsilon$。
\begin{solution}
对每个 $b_i\neq0$ 的 $A_i$（Borel 集，测度有限），由第~\ref{lec:extension-of-measures} 讲作业 C7) 取有限个左开右闭区间的并 $B_i$，使 $\lambda(A_i\triangle B_i)<\frac\varepsilon{m\abs{b_i}}$。令 $q=\sum_ib_i\one_{B_i}$。$\abs{\one_{A_i}-\one_{B_i}}=\one_{A_i\triangle B_i}$，所以
\[
\int\abs{\varphi-q}\dd x\leqslant\sum_i\abs{b_i}\lambda(A_i\triangle B_i)<\varepsilon .
\]
把所有区间的端点从小到大排成 $a_0<a_1<\cdots<a_k$，$q$ 在每个 $(a_{j-1},a_j]$ 上是常数 $c_j$，在 $(a_0,a_k]$ 之外为零，所以 $q=\sum_jc_j\one_{(a_{j-1},a_j]}$。由 F1) 与三角不等式，$\int\abs{g-q}<2\varepsilon$。
\end{solution}
\item 对 $q=\one_{(\alpha,\beta]}$ 直接计算 $\int q(x)\cos nx\dd x$，证明它趋于零。从而对一切阶梯函数 $q$ 成立。
\begin{solution}
由定理~\ref{thm:integration:riemann}，$\int_{(\alpha,\beta]}\cos nx\dd x=\frac{\sin n\beta-\sin n\alpha}n$，绝对值不超过 $\frac2n\to0$。阶梯函数是有限个这样的函数的线性组合，由线性性也趋于零。
\end{solution}
\item 由 $\abs{\int g\cos nx\dd x}\leqslant\int\abs{g-q}\dd x+\abs{\int q\cos nx\dd x}$ 完成证明。其中用到了积分的哪些性质？
\begin{solution}
$g\cos nx$ 与 $q\cos nx$ 可积（$\abs{\cos}\leqslant1$，由单调性）。由线性性，$\int g\cos nx=\int(g-q)\cos nx+\int q\cos nx$；由 $\abs{\int h}\leqslant\int\abs h$ 与单调性（$\abs{(g-q)\cos nx}\leqslant\abs{g-q}$），第一项的绝对值不超过 $\int\abs{g-q}<2\varepsilon$。由 F3)，$\limsup_n\abs{\int g\cos nx}\leqslant2\varepsilon$，$\varepsilon$ 任意。

用到的性质：线性性、单调性、$\abs{\int h}\leqslant\int\abs h$（定理~\ref{thm:integration:integrable}），以及 F1) 中的单调收敛定理。证明的思路是“先对简单的函数证明，再用逼近”：可积函数可以用阶梯函数在 $\int\abs{\cdot}$ 的意义下任意逼近。
\end{solution}
\item 证明同样有 $\int g(x)\sin nx\dd x\to0$。又设 $g$ 只在 $[a,b]$ 上可积，证明 $\int_a^bg(x)\cos nx\dd x\to0$。
\begin{solution}
F3) 中换成 $\int_{(\alpha,\beta]}\sin nx\dd x=\frac{\cos n\alpha-\cos n\beta}n\to0$，其余不变。第二句：对 $g\one_{[a,b]}$（它在 $\R$ 上可积）用已证的结论。
\end{solution}
\item （用连续函数逼近）在 F2) 中，把每个 $c_j\one_{(a_{j-1},a_j]}$ 换成连续函数 $r_j$：在 $(a_{j-1},a_j)$ 之外为 $0$，在 $[a_{j-1}+\delta,a_j-\delta]$ 上为 $c_j$，其余部分线性。证明 $\delta$ 充分小时 $r=\sum_jr_j$ 满足 $\int\abs{q-r}\dd x<\varepsilon$。由此，可积函数可以用具有紧支撑的连续函数在 $\int\abs{\cdot}$ 的意义下任意逼近。
\begin{solution}
取 $\delta<\frac12\min_j(a_j-a_{j-1})$。$c_j\one_{(a_{j-1},a_j]}$ 与 $r_j$ 只在 $(a_{j-1},a_{j-1}+\delta)\cup(a_j-\delta,a_j]$ 上不同，这里差的绝对值不超过 $\abs{c_j}$，所以 $\int\abs{c_j\one_{(a_{j-1},a_j]}-r_j}\leqslant2\delta\abs{c_j}$。求和，$\int\abs{q-r}\leqslant2\delta\sum_j\abs{c_j}<\varepsilon$（$\delta$ 充分小）。$r$ 连续，在 $[a_0,a_k]$ 之外为零。结合 F2)，$\int\abs{g-r}<3\varepsilon$。

这个逼近在分析中十分常用。比如要证明某个关于可积函数的结论，可以先对连续函数（甚至光滑函数）证明，再用逼近推广。
\end{solution}
\end{questions}
\end{probleme}
